\documentclass[nojss,noheadings]{jss}

\usepackage{amsmath,amssymb,amsfonts}
\usepackage{bm}
\usepackage{booktabs}
\usepackage{tabularx}
\usepackage{enumitem}
\usepackage{placeins}
\raggedbottom

\newcommand{\R}{\mathbb{R}}
\newcommand{\Rp}{\mathbb{R}_{+}}
\newcommand{\Ind}{\mathbb{I}}
\newcommand{\dd}{\mathrm{d}}
\renewcommand{\E}{\mathbb{E}}
\newcommand{\Var}{\operatorname{Var}}
\newcommand{\Cov}{\operatorname{Cov}}
\newcommand{\tr}{\operatorname{tr}}
\newcommand{\diag}{\operatorname{diag}}
\newcommand{\logit}{\operatorname{logit}}
\newcommand{\expit}{\operatorname{expit}}
\newcommand{\vect}[1]{\bm{#1}}
\newcommand{\mat}[1]{\mathbf{#1}}
\newcommand{\N}{\mathcal{N}}
\newcommand{\TNp}{\mathcal{N}^{+}}
\newcommand{\IG}{\operatorname{IG}}
\newcommand{\GIG}{\operatorname{GIG}}
\newcommand{\Exp}{\operatorname{Exp}}
\newcommand{\AL}{\operatorname{AL}}
\newcommand{\exAL}{\operatorname{exAL}}
\newcommand{\KL}{\operatorname{KL}}
\newcommand{\CRPS}{\operatorname{CRPS}}
\newcommand{\PPLC}{\operatorname{PPLC}}
\newcommand{\elbo}{\mathcal{L}}
\newcommand{\pigamma}{\pi_\gamma}
\newcommand{\TableStyle}{\footnotesize\setlength{\tabcolsep}{4pt}\renewcommand{\arraystretch}{1.08}}
\newcounter{suppalgorithm}
\renewcommand{\thesuppalgorithm}{S\arabic{suppalgorithm}}
\newenvironment{suppalgorithm}[1]{%
  \refstepcounter{suppalgorithm}%
  \par\medskip\noindent\textbf{Algorithm~\thesuppalgorithm. #1.}\par\nobreak
}{\par\medskip}
\renewcommand{\theequation}{S\arabic{equation}}
\renewcommand{\thetable}{S\arabic{table}}
\renewcommand{\thesection}{S\arabic{section}}
\renewcommand{\thesubsection}{S\arabic{section}.\arabic{subsection}}
\renewcommand{\thesubsubsection}{S\arabic{section}.\arabic{subsection}.\arabic{subsubsection}}

\author{Antonio Aguirre \\ University of California Santa Cruz \And
        Raquel Barata \\ University of California Santa Cruz \AND
        Raquel Prado \\ University of California Santa Cruz \And
        Bruno Sans\'{o} \\ University of California Santa Cruz}
\title{Supplementary Material for ``\pkg{exdqlm}: An \proglang{R} Package for Estimation and Analysis of Flexible Dynamic Quantile Linear Models''}

\Plainauthor{Antonio Aguirre, Raquel Barata, Raquel Prado, Bruno Sanso}
\Plaintitle{Supplementary Material for exdqlm: An R Package for Estimation and Analysis of Flexible Dynamic Quantile Linear Models}
\Shorttitle{Supplementary material for \pkg{exdqlm}}

\Abstract{This supplement derives the posterior targets, full conditional
distributions, variational updates, Laplace--delta approximations, evidence
lower bound (ELBO) terms, and numerical blocks used by the algorithms in
\pkg{exdqlm}. It covers state-space DQLMs under the asymmetric Laplace (AL)
likelihood and exDQLMs under the extended asymmetric Laplace (exAL) likelihood,
transfer-function state augmentation, static AL and exAL regression, ridge and
Nishimura--Suchard regularized horseshoe
shrinkage priors, forecasting, diagnostics, and posterior-predictive synthesis.
The notation follows the main article, and each algorithmic block is linked to
the corresponding exported package function.}

\Keywords{Bayesian quantile regression, dynamic linear models, asymmetric
Laplace likelihood, extended asymmetric Laplace likelihood, MCMC,
Laplace--delta variational Bayes}
\Plainkeywords{Bayesian quantile regression, dynamic linear models, asymmetric
Laplace likelihood, extended asymmetric Laplace likelihood, MCMC,
Laplace-delta variational Bayes}

\Address{
  Antonio Aguirre and Raquel Barata\\
  Department of Statistics\\
  University of California Santa Cruz\\
  Santa Cruz, CA 95064\\
  E-mail: \email{jaguir26@ucsc.edu}, \email{raquel.a.barata@gmail.com}\\
}

\begin{document}

\section{Scope and relation to the main article}
\label{sec:supp_scope}

The main article describes the user-facing workflow for \pkg{exdqlm}: dynamic
exDQLM and DQLM fitting, transfer-function models, static exAL regression,
forecasting, diagnostics, and posterior-predictive synthesis. This supplement
gives the corresponding posterior targets, full conditional distributions,
Markov chain Monte Carlo (MCMC) and variational Bayes (VB) updates,
Laplace--delta variational Bayes (LDVB) approximations, ELBO components, and
numerical routines used by the package. Readers can use this document to
verify the posterior targets, update equations, and diagnostic quantities
behind the exported functions. The derivations build on the AL
mixture representation of \citet{kozumi2011gibbs}, dynamic quantile linear
models \citep{gonccalves2017dynamic}, the exAL/exDQLM formulation
\citep{yan2025family,me}, Gaussian state-space filtering, smoothing, and
forward-filtering backward-sampling (FFBS) recursions described in
Section~\ref{sec:gaussian_state_update},
nonconjugate variational approximations \citep{wang2013nonconjugatevb}, and
the \code{rhs\_ns} prior \citep{nishimura2023shrunken}.

Table~\ref{tab:supp_roadmap} gives the main navigation map, linking each
article topic to the corresponding supplement sections and exported package
routines. The model sections derive the fitting targets and updates, and
Section~\ref{sec:numerical_blocks} collects reusable numerical identities.

The notation follows the main article. In particular, $\mat{F}_t$ denotes the
observation vector, $\mat{G}_t$ the state evolution matrix, $\mat{W}_t$ the
evolution covariance, $\vect{\theta}_t$ the dynamic state, $p_0$ the target
quantile level, and $\gamma$ the exAL skewness parameter. Dimensions are stated
when each object is introduced.

\begin{table}[!htb]
\centering
\begingroup
\TableStyle
\begin{tabularx}{\textwidth}{@{}>{\raggedright\arraybackslash}p{0.22\textwidth}
>{\raggedright\arraybackslash}p{0.18\textwidth}
>{\raggedright\arraybackslash}p{0.15\textwidth}
>{\raggedright\arraybackslash}X@{}}
\toprule
Article topic or workflow & Main article location & Supplement location &
Exported functions or computational block \\
\midrule
exAL distribution and gamma bounds &
Section 2.1 &
S2 &
\code{dexal()}, \code{pexal()}, \code{qexal()}, \code{rexal()},
\code{get\_gamma\_bounds()}. \\
Dynamic DQLM and exDQLM fitting &
Sections 2.1--2.2, Examples 1--2 &
S3--S4 &
\code{exdqlmMCMC()}, \code{exdqlmLDVB()}; dynamic fitting algorithms. \\
Gaussian state updates &
Section 2.2 &
S2.2 &
Kalman filtering, FFBS backward sampling, and RTS smoothing. \\
Discount-factor evolution &
Section 2.3 &
S2.2, S3.1, S4.1 &
State-evolution covariance construction used by the dynamic routines. \\
Transfer-function exDQLM &
Section 2.4, Example 3 &
S5 &
\code{exdqlmTransferMCMC()}, \code{exdqlmTransferLDVB()}. \\
Static AL and exAL regression &
Section 2.5, Example 4 &
S6--S7 &
\code{exalStaticMCMC()}, \code{exalStaticLDVB()}. \\
Forecasting &
Section 3.2 &
S8.1 &
\code{exdqlmForecast()}; forecast recursion and predictive draws. \\
Diagnostics &
Section 3.1, Examples 2--3 &
S8.2 &
\code{exdqlmDiagnostics()}, \code{exalStaticDiagnostics()}; calibration and
predictive-score summaries. \\
Posterior-predictive synthesis &
Example 1 &
S8.3 &
\code{quantileSynthesis()}; Algorithm~\ref{alg:predictive-synthesis}. \\
Reusable numerical blocks &
Section 2.2 &
S9 &
GIG moments, entropy terms, slice sampling, and Laplace--delta approximation. \\
\bottomrule
\end{tabularx}
\endgroup
\caption{Roadmap connecting the main article, supplement derivations, and
exported \pkg{exdqlm} routines or computational blocks.}
\label{tab:supp_roadmap}
\end{table}

Some routines temporarily hold or schedule selected blocks during early
iterations for numerical stability. These controls affect update timing only;
they do not change the posterior target, full conditional distributions, or
coordinate optima documented in the sections listed in
Table~\ref{tab:supp_roadmap}.

\section{Notation and distributional conventions}
\label{sec:supp_distributional_conventions}

Let $p_0\in(0,1)$ be the target quantile. The response is scalar. In dynamic
models, observations are indexed by $t=1,\ldots,T$; in static regression,
observations are indexed by $i=1,\ldots,n$. For any positive definite matrix
$\mat{M}$, $|\mat{M}|$ denotes its determinant.

\subsection{Core notation}

The symbols $p_\gamma$, $A_\gamma$, $B_\gamma$, and $C_\gamma$ are deterministic
functions of $(p_0,\gamma)$, whereas $\pigamma(\gamma)$ denotes the prior
density for $\gamma$ on $(L,U)$. Generic densities are denoted by $p(\cdot)$
only when no confusion with the exAL quantile map can arise. In variational
sections, expectations without an explicit subscript are taken under the
current mean-field distribution $q$. The symbol $\xi$ is reserved for the
\code{rhs\_ns} half-Cauchy auxiliary variable; the transformed coordinate for $\gamma$
in Laplace--delta calculations is denoted by $z_\gamma$.

\begin{table}[!htb]
\centering
\begingroup
\TableStyle
\begin{tabularx}{\textwidth}{@{}lX@{}}
\toprule
Symbol & Meaning \\
\midrule
$p_0$ & Target quantile level. \\
$p_\gamma=p(p_0,\gamma)$ & Internal exAL quantile map defined in the main article. \\
$A_{p_0},B_{p_0}$ & AL constants $A(p_0)$ and $B(p_0)$. \\
$A_\gamma,B_\gamma,C_\gamma$ & exAL constants evaluated at $p_\gamma$. \\
$L,U$ & Lower and upper admissible bounds for $\gamma$. \\
$\pigamma(\gamma)$ & Prior density for $\gamma$ on $(L,U)$. \\
$\vect{\theta}_t,\mat{F}_t,\mat{G}_t,\mat{W}_t$ & Dynamic state, observation vector, evolution matrix, and evolution covariance. \\
$d_t,Q_t$ & Observation offset and variance in the generic Gaussian state update. \\
$v_t,s_t$ & exAL latent variables in the dynamic model. \\
$\vect{\beta},\mat{X}$ & Static regression coefficients and design matrix. \\
$v_i,s_i$ & exAL latent variables in the static model. \\
$\mathcal{A}$ & Active coefficient set receiving \code{rhs\_ns} shrinkage. \\
$\elbo(q)$ & Evidence lower bound. \\
\bottomrule
\end{tabularx}
\endgroup
\caption{Core notation used throughout the supplement. The shorthand
$p_\gamma$, $A_\gamma$, $B_\gamma$, and $C_\gamma$ refers to functions of
$(p_0,\gamma)$ defined in the main article.}
\label{tab:supp_notation}
\end{table}

The latent variables $v_t$, $s_t$, $v_i$, and $s_i$ are introduced for
posterior computation and are updated internally by the fitting routines; they
are not user-supplied inputs. Practitioners specify the response, target
quantile, model components, priors, and inference controls described in the
main article.

The distributional parameterizations are also fixed throughout:
$\N(m,V)$ uses variance $V$;
\begin{equation}
x\sim\IG(a,b)
\quad\Longleftrightarrow\quad
p(x)=\frac{b^a}{\Gamma(a)}x^{-a-1}\exp(-b/x),\qquad x>0,
\end{equation}
equivalently $x^{-1}\sim\operatorname{Gamma}(a,b)$ under a shape-rate gamma
parameterization. The exponential distribution $\Exp(\text{rate}=1/\sigma)$
has density $\sigma^{-1}\exp(-x/\sigma)$, matching the main article's
mean-$\sigma$ convention for the AL and exAL mixture variables. Finally,
$\TNp(m,V)$ denotes $\N(m,V)$ truncated to $(0,\infty)$. The generalized
inverse Gaussian (GIG) parameterization is given in \eqref{eq:gig_parameterization}.
Throughout, $\phi(\cdot)$ and $\Phi(\cdot)$ denote the standard normal density
and distribution function.

The asymmetric Laplace (AL) mixture representation used by the package is
\begin{align}
v &\mid \sigma \sim \Exp(\text{rate}=1/\sigma), \qquad v>0, \\
y \mid \eta,\sigma,v
&\sim \N\{ \eta + A(p_0)v,\ \sigma B(p_0)v \},
\end{align}
where $\eta$ is the quantile-location parameter and
\begin{equation}
A(p)=\frac{1-2p}{p(1-p)}, \qquad B(p)=\frac{2}{p(1-p)}.
\label{eq:al_AB}
\end{equation}
The AL model is the $\gamma=0$ special case of the exAL formulation below.

The exAL distribution is represented through an additional latent variable
$s\sim\TNp(0,1)$, a standard normal truncated to $(0,\infty)$. For
$\gamma\in(L,U)$, define
\begin{align}
g(\gamma) &= 2\Phi(-|\gamma|)\exp(\gamma^2/2), \\
p_\gamma &= \Ind(\gamma<0) + \frac{p_0-\Ind(\gamma<0)}{g(\gamma)}, \\
A_\gamma &= A(p_\gamma), \qquad
B_\gamma = B(p_\gamma), \qquad
C_\gamma = \{\Ind(\gamma>0)-p_\gamma\}^{-1}.
\label{eq:exal_ABC}
\end{align}
The endpoints $L<0<U$ are the roots $g(L)=1-p_0$ and $g(U)=p_0$. Conditional on
$(v,s,\sigma,\gamma)$,
\begin{equation}
y\mid \eta,\sigma,\gamma,v,s
\sim
\N\left(\eta + C_\gamma\sigma|\gamma|s + A_\gamma v,\ 
\sigma B_\gamma v \right).
\label{eq:exal_aug}
\end{equation}
The package uses an inverse-gamma prior $\sigma\sim\IG(a_\sigma,b_\sigma)$
and a proper truncated Student-$t$ prior $\pigamma(\gamma)$ for $\gamma$ on
$(L,U)$. If $m_\gamma$, $s_\gamma$, and $\nu_\gamma$ denote its location,
scale, and degrees of freedom, then
\begin{equation}
\pigamma(\gamma)=
\frac{
s_\gamma^{-1}t_{\nu_\gamma}\{(\gamma-m_\gamma)/s_\gamma\}
}{
T_{\nu_\gamma}\{(U-m_\gamma)/s_\gamma\}
-T_{\nu_\gamma}\{(L-m_\gamma)/s_\gamma\}
}\Ind(L<\gamma<U),
\label{eq:gamma_prior}
\end{equation}
where $t_\nu(\cdot)$ and $T_\nu(\cdot)$ are the standard Student-$t$ density
and distribution function.

The generalized inverse Gaussian distribution is parameterized as
\begin{equation}
x\sim\GIG(\lambda,\chi,\psi_{\mathrm{GIG}})
\quad\Longleftrightarrow\quad
p(x)\propto x^{\lambda-1}
\exp\left\{-\frac{1}{2}\left(\frac{\chi}{x}+\psi_{\mathrm{GIG}} x\right)\right\},
\qquad x>0.
\label{eq:gig_parameterization}
\end{equation}
For $r\in\R$,
\begin{equation}
\E(x^r)=
\left(\frac{\chi}{\psi_{\mathrm{GIG}}}\right)^{r/2}
\frac{K_{\lambda+r}(\sqrt{\chi\psi_{\mathrm{GIG}}})}
{K_{\lambda}(\sqrt{\chi\psi_{\mathrm{GIG}}})},
\label{eq:gig_moments}
\end{equation}
where $K_\nu(\cdot)$ is the modified Bessel function of the second kind.

\subsection{Gaussian state updates for dynamic models}
\label{sec:gaussian_state_update}

All dynamic MCMC and VB updates for the state path reduce to the same
conditionally Gaussian state-space calculation. The package uses the
observation-offset representation
\begin{equation}
y_t=\mat{F}_t^\top\vect{\theta}_t+d_t+e_t,\qquad
e_t\sim\N(0,Q_t),
\label{eq:generic_gaussian_obs}
\end{equation}
where $d_t$ is the scalar observation offset and $Q_t>0$ is the scalar
conditional observation variance. In the implementation these quantities are
stored as the extra observation offset and variance, respectively. Equivalently,
with $z_t=y_t-d_t$,
\[
z_t=\mat{F}_t^\top\vect{\theta}_t+e_t,\qquad e_t\sim\N(0,Q_t).
\]
Together with
\[
\vect{\theta}_0\sim\N(\vect{m}_0,\mat{C}_0),\qquad
\vect{\theta}_t\mid\vect{\theta}_{t-1}
\sim
\N(\mat{G}_t\vect{\theta}_{t-1},\mat{W}_t),
\]
this defines the Gaussian state update used below.

Initialize $\vect{m}_{0|0}=\vect{m}_0$ and
$\mat{C}_{0|0}=\mat{C}_0$. For $t=1,\ldots,T$, the Kalman prediction and
filtering recursions are
\begin{align}
\vect{a}_t
&=\mat{G}_t\vect{m}_{t-1|t-1},&
\mat{R}_t
&=\mat{G}_t\mat{C}_{t-1|t-1}\mat{G}_t^\top+\mat{W}_t,
\label{eq:generic_kalman_prediction}\\
\mu_t^y
&=d_t+\mat{F}_t^\top\vect{a}_t,&
S_t
&=\mat{F}_t^\top\mat{R}_t\mat{F}_t+Q_t,
\label{eq:generic_kalman_forecast}\\
\delta_t^y
&=y_t-\mu_t^y,&
\vect{k}_t
&=\mat{R}_t\mat{F}_tS_t^{-1},
\label{eq:generic_kalman_gain}\\
\vect{m}_{t|t}
&=\vect{a}_t+\vect{k}_t\delta_t^y,&
\mat{C}_{t|t}
&=\mat{R}_t-\vect{k}_tS_t\vect{k}_t^\top.
\label{eq:generic_kalman_filter}
\end{align}
The same forward pass gives the Gaussian auxiliary normalizing constant
\begin{equation}
\log p_{\mathrm{G}}(\vect{y})
=
\sum_{t=1}^T
\log \N\{y_t\mid d_t+\mat{F}_t^\top\vect{a}_t,S_t\},
\label{eq:generic_kalman_loglik}
\end{equation}
which is used in the ELBO state-block identities below.

For MCMC, the state path is sampled by forward-filtering backward-sampling
(FFBS) \citep{carter1994gibbs,fruhwirth1994data}. After the forward filter,
draw
\[
\vect{\theta}_T\sim\N(\vect{m}_{T|T},\mat{C}_{T|T}).
\]
For $t=T-1,\ldots,0$, define
\begin{equation}
\mat{J}_t
=
\mat{C}_{t|t}\mat{G}_{t+1}^\top\mat{R}_{t+1}^{-1}.
\label{eq:generic_backward_gain}
\end{equation}
Then sample
\begin{equation}
\vect{\theta}_t\mid\vect{\theta}_{t+1},\vect{y}
\sim
\N(\vect{h}_t,\mat{H}_t),
\label{eq:generic_backward_sampling}
\end{equation}
where
\begin{align}
\vect{h}_t
&=
\vect{m}_{t|t}
+
\mat{J}_t(\vect{\theta}_{t+1}-\vect{a}_{t+1}),\\
\mat{H}_t
&=
\mat{C}_{t|t}
-
\mat{J}_t\mat{R}_{t+1}\mat{J}_t^\top.
\label{eq:generic_backward_sampling_moments}
\end{align}

For VB and LDVB, the state factor $q(\vect{\theta}_{0:T})$ is updated
deterministically by the Rauch--Tung--Striebel (RTS) backward smoother applied to the
same Gaussian system. Set $\vect{m}_{T|T}$ and $\mat{C}_{T|T}$ equal to the
final filtered moments. For $t=T-1,\ldots,0$,
\begin{align}
\vect{m}_{t|T}
&=
\vect{m}_{t|t}
+
\mat{J}_t(\vect{m}_{t+1|T}-\vect{a}_{t+1}),\\
\mat{C}_{t|T}
&=
\mat{C}_{t|t}
+
\mat{J}_t(\mat{C}_{t+1|T}-\mat{R}_{t+1})\mat{J}_t^\top.
\label{eq:generic_rts_smoother}
\end{align}
Thus MCMC uses stochastic backward sampling from the Gaussian conditional state
posterior, whereas VB/LDVB uses deterministic smoothing to update the Gaussian
factor $q(\vect{\theta}_{0:T})$. Posterior samples returned by the package from
VB/LDVB fits are generated after convergence from the fitted variational
factors; this sampling step is separate from the coordinate update.

The discount-factor construction changes how $\mat{W}_t$, or equivalently
$\mat{R}_t$, is formed during the prediction step
\citep[Section~6.3]{west1997bayesian}. Let
\[
\mat{P}_t
=
\mat{G}_t\mat{C}_{t-1|t-1}\mat{G}_t^\top
\]
denote the propagated covariance before discount inflation. For a partition of
the state vector into blocks $\mathcal{B}_1,\ldots,\mathcal{B}_h$, with
discount factors $\delta_1,\ldots,\delta_h$, the package constructs a
blockwise inflation matrix $\mat{D}$ with entries
\[
D_{ij}
=
\begin{cases}
(1-\delta_b)/\delta_b, & i,j\in\mathcal{B}_b,\\
0, & i\in\mathcal{B}_b,\ j\in\mathcal{B}_{b'},\ b\neq b',
\end{cases}
\]
and uses
\[
\mat{W}_t=\mat{D}\circ\mat{P}_t,\qquad
\mat{R}_t=\mat{P}_t+\mat{W}_t,
\]
where $\circ$ denotes elementwise multiplication. Thus within-block entries of
the propagated covariance are inflated by $1/\delta_b$, while cross-block
entries of $\mat{R}_t$ remain those in $\mat{P}_t$. If all states share a single
discount factor $\delta$, this reduces to the scalar-discount special case
$\mat{R}_t=\delta^{-1}\mat{P}_t$ and
$\mat{W}_t=\{(1-\delta)/\delta\}\mat{P}_t$. Once $\mat{R}_t$ is formed, the
filtering, backward-sampling, and smoothing equations above are unchanged.

This subsection is model-agnostic. To use it for any dynamic model below,
select the appropriate row of Table~\ref{tab:gaussian_state_update_map}, plug
that row's offset $d_t$ and variance $Q_t$ into
\eqref{eq:generic_gaussian_obs}, and then apply the same filtering and backward
recursions above. The variational rows are obtained by completing the square in
the expected Gaussian log likelihood.
\begin{table}[!htb]
\centering
\begingroup
\TableStyle
\begin{tabularx}{\textwidth}{@{}>{\raggedright\arraybackslash}p{0.20\textwidth}
>{\raggedright\arraybackslash}p{0.25\textwidth}
>{\raggedright\arraybackslash}p{0.19\textwidth}
>{\raggedright\arraybackslash}X@{}}
\toprule
Dynamic update & Offset $d_t$ & Variance $Q_t$ & State update \\
\midrule
DQLM MCMC &
$A_{p_0}v_t$ &
$\sigma B_{p_0}v_t$ &
Kalman filter and FFBS backward sampling. \\
exDQLM MCMC &
$A_\gamma v_t+C_\gamma\sigma|\gamma|s_t$ &
$\sigma B_\gamma v_t$ &
Kalman filter and FFBS backward sampling. \\
DQLM VB &
$A_{p_0}/m_{1/v,t}$ &
$B_{p_0}/\{\kappa m_{1/v,t}\}$ &
Kalman filter and deterministic RTS smoothing. \\
exDQLM LDVB &
$\{\kappa_2+\kappa_4\bar s_t m_{1/v,t}\}/\phi_t$ &
$\phi_t^{-1}$, $\phi_t=\kappa_1m_{1/v,t}$ &
Kalman filter and deterministic RTS smoothing. \\
Transfer-function exDQLM &
Same as the selected base dynamic update &
Same as the selected base dynamic update &
Apply the same update after replacing
$(\mat{F}_t,\mat{G}_t,\mat{W}_t,\vect{\theta}_t)$ by the augmented quantities
in Section~\ref{sec:transfer_function}. \\
\bottomrule
\end{tabularx}
\endgroup
\caption{Offset and variance representation used by the generic Gaussian state
update in \eqref{eq:generic_gaussian_obs}. Here
$\kappa=\E(\sigma^{-1})$ in the DQLM VB row,
$m_{1/v,t}=\E(v_t^{-1})$, $\bar s_t=\E(s_t)$, and the $\kappa_j$ moments are
defined in \eqref{eq:dyn_exdqlm_kappa}.}
\label{tab:gaussian_state_update_map}
\end{table}
\FloatBarrier

\section{Dynamic DQLM under the AL likelihood}
\label{sec:dyn_dqlm}

\subsection{Model hierarchy and posterior target}

Let $\vect{\theta}_t\in\R^q$ be the latent state, $\mat{F}_t\in\R^q$ the
observation vector, $\mat{G}_t\in\R^{q\times q}$ the evolution matrix, and
$\mat{W}_t\in\R^{q\times q}$ a positive semidefinite evolution covariance. The
dynamic DQLM is
\begin{align}
\vect{\theta}_0 &\sim \N(\vect{m}_0,\mat{C}_0), \qquad
  \vect{m}_0\in\R^q,\ \mat{C}_0\in\R^{q\times q}, \\
\vect{\theta}_t\mid\vect{\theta}_{t-1}
&\sim \N(\mat{G}_t\vect{\theta}_{t-1},\mat{W}_t), \\
v_t\mid\sigma &\sim \Exp(\text{rate}=1/\sigma), \\
y_t\mid\vect{\theta}_t,\sigma,v_t
&\sim \N(\mat{F}_t^\top\vect{\theta}_t + A_{p_0} v_t,\ \sigma B_{p_0} v_t),
\label{eq:dyn_dqlm_hierarchy}
\end{align}
where $A_{p_0}=A(p_0)$ and $B_{p_0}=B(p_0)$. With
$\vect{v}=(v_1,\ldots,v_T)^\top$ and
$\vect{\theta}_{0:T}=(\vect{\theta}_0,\ldots,\vect{\theta}_T)$, the posterior
target is
\begin{align}
p(\vect{\theta}_{0:T},\sigma,\vect{v}\mid\vect{y})
&\propto
p(\sigma)\,p(\vect{\theta}_0)
\prod_{t=1}^T p(\vect{\theta}_t\mid\vect{\theta}_{t-1})
\prod_{t=1}^T p(v_t\mid\sigma)
\prod_{t=1}^T p(y_t\mid\vect{\theta}_t,\sigma,v_t).
\label{eq:dyn_dqlm_joint}
\end{align}
Equivalently, conditional on $(\sigma,\vect{v})$, the pseudo-observation
$\tilde y_t=y_t-A_{p_0}v_t$ satisfies
\begin{equation}
\tilde y_t=\mat{F}_t^\top\vect{\theta}_t+\epsilon_t,\qquad
\epsilon_t\sim\N(0,\sigma B_{p_0}v_t).
\label{eq:dyn_dqlm_pseudo}
\end{equation}
Thus the conditional state posterior is the Gaussian state-space posterior
defined in Section~\ref{sec:gaussian_state_update}, with the DQLM MCMC offset
and variance given in Table~\ref{tab:gaussian_state_update_map}. The package
constructs the evolution covariance through the discount-factor state-space
rules described in the main article; once the prediction covariance is formed,
the filtering and backward updates are those in
Section~\ref{sec:gaussian_state_update}.

\subsection{MCMC updates}

One Gibbs sweep for the dynamic DQLM uses the following full conditionals.

\paragraph{State path.}
Given $(\sigma,\vect{v})$, sample $\vect{\theta}_{0:T}$ by the FFBS recursion
in Section~\ref{sec:gaussian_state_update}, using
$d_t=A_{p_0}v_t$ and $Q_t=\sigma B_{p_0}v_t$.

\paragraph{Scale.}
Let $r_t=y_t-\mat{F}_t^\top\vect{\theta}_t$. Then
\begin{equation}
\sigma\mid \vect{\theta}_{0:T},\vect{v},\vect{y}
\sim \IG\left(
a_\sigma+\frac{3T}{2},\
b_\sigma+\sum_{t=1}^T v_t+
\sum_{t=1}^T\frac{(r_t-A_{p_0}v_t)^2}{2B_{p_0}v_t}
\right).
\label{eq:dyn_dqlm_sigma_cond}
\end{equation}

\paragraph{Latent AL variables.}
For $t=1,\ldots,T$,
\begin{equation}
v_t\mid\vect{\theta}_t,\sigma,y_t
\sim
\GIG\left(
\frac{1}{2},\
\frac{r_t^2}{\sigma B_{p_0}},\
\frac{A_{p_0}^2}{\sigma B_{p_0}}+\frac{2}{\sigma}
\right).
\label{eq:dyn_dqlm_v_cond}
\end{equation}

\begin{suppalgorithm}{Dynamic DQLM MCMC}
\label{alg:dyn-dqlm-mcmc}
The MCMC kernel leaves the augmented posterior in
\eqref{eq:dyn_dqlm_joint} invariant; posterior summaries are Monte Carlo
estimates from the retained finite chain.
\par\noindent\emph{Input:} observations $\vect{y}$, model matrices
$(\mat{F}_t,\mat{G}_t,\mat{W}_t)$, initial state prior, AL prior
hyperparameters, and MCMC controls $N_{\mathrm{burn}}$ and
$N_{\mathrm{mcmc}}$.
\par\noindent\emph{Output:} retained posterior draws for
$(\vect{\theta}_{0:T},\vect{v},\sigma)$ and filtered/smoothed state summaries.
\begin{enumerate}[label=(\arabic*)]
\item Initialize $\vect{\theta}_{0:T}$, $\vect{v}$, and $\sigma$.
\item For $m=1,\ldots,N_{\mathrm{burn}}+N_{\mathrm{mcmc}}$:
  \begin{enumerate}[label=(\alph*)]
  \item sample $\vect{\theta}_{0:T}$ by FFBS using
  Section~\ref{sec:gaussian_state_update} and the DQLM MCMC row of
  Table~\ref{tab:gaussian_state_update_map};
  \item sample $\sigma$ from \eqref{eq:dyn_dqlm_sigma_cond};
  \item sample each $v_t$ from \eqref{eq:dyn_dqlm_v_cond};
  \item if $m>N_{\mathrm{burn}}$, retain the current draw.
  \end{enumerate}
\item Return the retained draws and state summaries.
\end{enumerate}
\end{suppalgorithm}

\subsection{Mean-field VB and ELBO}

The dynamic DQLM VB approximation uses
\begin{equation}
q(\vect{\theta}_{0:T},\sigma,\vect{v})
=q(\vect{\theta}_{0:T})q(\sigma)\prod_{t=1}^Tq(v_t),
\label{eq:dyn_dqlm_mf}
\end{equation}
with $q(\sigma)=\IG(a_q,b_q)$ and
$q(v_t)=\GIG(1/2,\chi_{v,t},\psi_{\mathrm{GIG},t})$. Define
\begin{equation}
\kappa=\E_q(\sigma^{-1})=\frac{a_q}{b_q},\qquad
m_{1/v,t}=\E_q(v_t^{-1}),\qquad
\bar v_t=\E_q(v_t).
\end{equation}
The state factor $q(\vect{\theta}_{0:T})$ is a joint Gaussian smoothing
distribution, not a product over time. The latent-variable updates require
only its marginal means and variances.
The state factor is the Gaussian posterior of the auxiliary model
\begin{equation}
y_t^\star=\mat{F}_t^\top\vect{\theta}_t+\epsilon_t^\star,\qquad
\epsilon_t^\star\sim\N(0,R_t^\star),
\qquad
y_t^\star=y_t-\frac{A_{p_0}}{m_{1/v,t}},\quad
R_t^\star=\frac{B_{p_0}}{\kappa m_{1/v,t}}.
\label{eq:dyn_dqlm_vb_pseudo}
\end{equation}
Equivalently, this is the generic observation-offset form in
\eqref{eq:generic_gaussian_obs} with
$d_t=A_{p_0}/m_{1/v,t}$ and
$Q_t=B_{p_0}/\{\kappa m_{1/v,t}\}$.
After applying the deterministic smoother in
Section~\ref{sec:gaussian_state_update}, write
$q(\vect{\theta}_t)=\N(\vect{m}_{t|T},\mat{C}_{t|T})$ and define
\begin{align}
\bar r_t &= y_t-\mat{F}_t^\top\vect{m}_{t|T},\\
\overline{r_t^2}
&= \bar r_t^2+\mat{F}_t^\top\mat{C}_{t|T}\mat{F}_t.
\end{align}
The remaining coordinate-ascent variational inference (CAVI) updates are
\begin{align}
q(v_t)&=\GIG\left(
\frac{1}{2},\
\frac{\kappa}{B_{p_0}}\overline{r_t^2},\
\kappa\left(2+\frac{A_{p_0}^2}{B_{p_0}}\right)\right), \\
q(\sigma)&=\IG(a_q,b_q), \qquad
a_q=a_\sigma+\frac{3T}{2},\\
b_q&=b_\sigma+\sum_{t=1}^T\bar v_t+
\frac{1}{2B_{p_0}}\sum_{t=1}^T
\left\{m_{1/v,t}\overline{r_t^2}-2A_{p_0}\bar r_t+A_{p_0}^2\bar v_t\right\}.
\label{eq:dyn_dqlm_vb_sigma}
\end{align}

\begin{suppalgorithm}{Dynamic DQLM VB}
\label{alg:dyn-dqlm-vb}
\par\noindent\emph{Input:} observations $\vect{y}$, model matrices
$(\mat{F}_t,\mat{G}_t,\mat{W}_t)$, initial state prior, AL prior
hyperparameters, convergence tolerance, and optional posterior-sampling
control \code{n.samp}.
\par\noindent\emph{Output:} fitted variational factors for
$(\vect{\theta}_{0:T},\vect{v},\sigma)$, convergence diagnostics, ELBO values
when requested, and optional approximate posterior draws.
\begin{enumerate}[label=(\arabic*)]
\item Initialize the variational parameters for $q(\vect{\theta}_{0:T})$,
$q(\sigma)$, and $q(v_t)$.
\item Repeat until the convergence criterion is met:
  \begin{enumerate}[label=(\alph*)]
  \item update $q(\vect{\theta}_{0:T})$ by the deterministic smoother in
  Section~\ref{sec:gaussian_state_update}, using the DQLM VB row of
  Table~\ref{tab:gaussian_state_update_map};
  \item update $q(v_t)$ for $t=1,\ldots,T$ using the GIG CAVI update above;
  \item update $q(\sigma)$ using \eqref{eq:dyn_dqlm_vb_sigma};
  \item monitor moment changes and, when ELBO monitoring is enabled, $\elbo(q)$.
  \end{enumerate}
\item Draw approximate posterior samples from the fitted variational factors
when requested by \code{n.samp}.
\item Return the fitted factors, diagnostics, and optional draws.
\end{enumerate}
\end{suppalgorithm}
In the package, this AL/DQLM path is selected in the dynamic fitting routines
with \code{dqlm.ind = TRUE}.

The ELBO is
\begin{align}
\elbo(q)
&=
\E_q\{\log p(\vect{y},\vect{\theta}_{0:T},\sigma,\vect{v})\}
-\E_q\{\log q(\vect{\theta}_{0:T},\sigma,\vect{v})\} \\
&=\elbo_\theta+\E_q\{\log p(\sigma)\}
+\sum_{t=1}^T\E_q\{\log p(v_t\mid\sigma)\}
+\sum_{t=1}^T\E_q\{\log p(y_t\mid\vect{\theta}_t,\sigma,v_t)\}\\
&\quad
+H\{q(\sigma)\}+\sum_{t=1}^TH\{q(v_t)\}.
\label{eq:dyn_dqlm_elbo}
\end{align}
The state-path term $\elbo_\theta$ is evaluated from the auxiliary Gaussian
state-space model in \eqref{eq:dyn_dqlm_vb_pseudo} through the Kalman marginal
likelihood identity
\begin{equation}
\elbo_\theta
=\log p(\vect{y}^{\star})
-\E_q\{\log p(\vect{y}^{\star}\mid\vect{\theta}_{0:T})\},
\label{eq:dyn_dqlm_elbo_theta}
\end{equation}
where $p(\vect{y}^{\star})$ is the marginal likelihood of the auxiliary
Gaussian state-space model. This identity is equivalent to
$\E_q\{\log p(\vect{\theta}_{0:T})\}+H\{q(\vect{\theta}_{0:T})\}$ for the
current state update.

The remaining terms in \eqref{eq:dyn_dqlm_elbo} are computable from the
current variational moments. Let
$\bar\ell_\sigma=\E(\log\sigma)$ and
$\bar\ell_{v,t}=\E(\log v_t)$. Then
\begin{align}
\E_q\{\log p(\sigma)\}
&=
a_\sigma\log b_\sigma-\log\Gamma(a_\sigma)
-(a_\sigma+1)\bar\ell_\sigma-b_\sigma\kappa,\\
\sum_{t=1}^T\E_q\{\log p(v_t\mid\sigma)\}
&=
-T\bar\ell_\sigma-\kappa\sum_{t=1}^T\bar v_t,\\
\sum_{t=1}^T\E_q\{\log p(y_t\mid\vect{\theta}_t,\sigma,v_t)\}
&=
-\frac{T}{2}\log(2\pi)-\frac{T}{2}\log B_{p_0}
-\frac{T}{2}\bar\ell_\sigma
-\frac{1}{2}\sum_{t=1}^T\bar\ell_{v,t}\nonumber\\
&\quad
-\frac{\kappa}{2B_{p_0}}\sum_{t=1}^T
\left\{m_{1/v,t}\overline{r_t^2}-2A_{p_0}\bar r_t+A_{p_0}^2\bar v_t\right\}.
\label{eq:dyn_dqlm_elbo_terms}
\end{align}
The entropy terms are given by \eqref{eq:entropy_ig} for $q(\sigma)$ and by
\eqref{eq:entropy_gig} for the $q(v_t)$ factors.

\section{Dynamic exDQLM under the exAL likelihood}
\label{sec:dyn_exdqlm}

\subsection{Model hierarchy and posterior target}

The dynamic exDQLM replaces the AL observation in
\eqref{eq:dyn_dqlm_hierarchy} with the exAL augmentation
\begin{align}
v_t\mid\sigma &\sim \Exp(\text{rate}=1/\sigma), \qquad
s_t\sim\TNp(0,1), \\
y_t\mid\vect{\theta}_t,\sigma,\gamma,v_t,s_t
&\sim
\N\left(
\mat{F}_t^\top\vect{\theta}_t+C_\gamma\sigma|\gamma|s_t+A_\gamma v_t,\
\sigma B_\gamma v_t
\right),
\label{eq:dyn_exdqlm_obs}
\end{align}
with the same state evolution as in Section~\ref{sec:dyn_dqlm}. The complete
posterior target is
\begin{align}
p(\vect{\theta}_{0:T},\vect{v},\vect{s},\sigma,\gamma\mid\vect{y})
&\propto
p(\vect{\theta}_0)\prod_{t=1}^Tp(\vect{\theta}_t\mid\vect{\theta}_{t-1})
\prod_{t=1}^Tp(y_t\mid\vect{\theta}_t,v_t,s_t,\sigma,\gamma) \nonumber\\
&\quad\times
\prod_{t=1}^Tp(v_t\mid\sigma)\prod_{t=1}^Tp(s_t)\,p(\sigma)\pigamma(\gamma).
\label{eq:dyn_exdqlm_joint}
\end{align}
Conditional on $(\vect{v},\vect{s},\sigma,\gamma)$,
\begin{equation}
\tilde y_t=y_t-A_\gamma v_t-C_\gamma\sigma|\gamma|s_t,\qquad
R^y_t=\sigma B_\gamma v_t,
\label{eq:dyn_exdqlm_pseudo}
\end{equation}
This is the generic Gaussian state update in
Section~\ref{sec:gaussian_state_update} with
$d_t=A_\gamma v_t+C_\gamma\sigma|\gamma|s_t$ and
$Q_t=\sigma B_\gamma v_t$.

\subsection{MCMC updates}

The default package path samples $\sigma$ from its exact conditional
given $\gamma$ and then updates $\gamma$ by bounded univariate slice sampling.
Joint random-walk alternatives are also implemented, but the following target
is the default exact-kernel path.

\paragraph{Latent $v_t$.}
Let
$r_t=y_t-\mat{F}_t^\top\vect{\theta}_t-C_\gamma\sigma|\gamma|s_t$. Then
\begin{equation}
v_t\mid\cdot
\sim
\GIG\left(
\frac{1}{2},\
\frac{r_t^2}{\sigma B_\gamma},\
\frac{A_\gamma^2}{\sigma B_\gamma}+\frac{2}{\sigma}
\right).
\label{eq:dyn_exdqlm_v_cond}
\end{equation}

\paragraph{Latent $s_t$.}
Let $y_t^\circ=y_t-\mat{F}_t^\top\vect{\theta}_t-A_\gamma v_t$ and
$D_\gamma=C_\gamma\sigma|\gamma|$. Then
\begin{equation}
s_t\mid\cdot\sim\TNp(\mu_{s,t},V_{s,t}),
\qquad
V_{s,t}=\left(1+\frac{D_\gamma^2}{\sigma B_\gamma v_t}\right)^{-1},
\qquad
\mu_{s,t}=V_{s,t}\frac{D_\gamma y_t^\circ}{\sigma B_\gamma v_t}.
\label{eq:dyn_exdqlm_s_cond}
\end{equation}

\paragraph{State path.}
Given $(\vect{v},\vect{s},\sigma,\gamma)$, sample
$\vect{\theta}_{0:T}$ by the FFBS recursion in
Section~\ref{sec:gaussian_state_update}, using the exDQLM MCMC row of
Table~\ref{tab:gaussian_state_update_map}.

\paragraph{Scale.}
For fixed $\gamma$, let $y_t^\circ=y_t-\mat{F}_t^\top\vect{\theta}_t-A_\gamma v_t$
and $c_{\gamma,t}=C_\gamma|\gamma|s_t$. Then
\begin{equation}
\sigma\mid\cdot
\sim
\GIG(\lambda_\sigma,\chi_\sigma,\psi_\sigma),
\label{eq:dyn_exdqlm_sigma_cond}
\end{equation}
where
\begin{align}
\lambda_\sigma &= -\left(a_\sigma+\frac{3T}{2}\right), \\
\chi_\sigma &= 2b_\sigma+2\sum_{t=1}^T v_t+
\sum_{t=1}^T\frac{(y_t^\circ)^2}{B_\gamma v_t},\\
\psi_\sigma &= \sum_{t=1}^T\frac{c_{\gamma,t}^2}{B_\gamma v_t}.
\end{align}

\paragraph{Skewness.}
The conditional kernel for $\gamma$ is
\begin{align}
\pi(\gamma\mid\cdot)
&\propto \pigamma(\gamma)\,
\prod_{t=1}^T(\sigma B_\gamma v_t)^{-1/2}
\exp\left[
-\sum_{t=1}^T
\frac{
\{y_t-\mat{F}_t^\top\vect{\theta}_t-A_\gamma v_t-C_\gamma\sigma|\gamma|s_t\}^2
}{2\sigma B_\gamma v_t}
\right],
\label{eq:dyn_exdqlm_gamma_kernel}
\end{align}
for $\gamma\in(L,U)$. The default sampler applies the bounded slice update in
Section~\ref{sec:slice} directly to this kernel.

\begin{suppalgorithm}{Dynamic exDQLM MCMC}
\label{alg:dyn-exdqlm-mcmc}
The default $(\sigma,\gamma)$ update consists of an exact GIG update for
$\sigma$ conditional on $\gamma$, followed by a bounded univariate slice update
for $\gamma$. Joint random-walk alternatives are implemented, but they are not
the default kernel summarized here.
\par\noindent\emph{Input:} observations $\vect{y}$, model matrices
$(\mat{F}_t,\mat{G}_t,\mat{W}_t)$, initial state prior, exAL prior
hyperparameters, admissible interval $(L,U)$ for $\gamma$, and MCMC controls
$N_{\mathrm{burn}}$ and $N_{\mathrm{mcmc}}$.
\par\noindent\emph{Output:} retained posterior draws for
$(\vect{\theta}_{0:T},\vect{v},\vect{s},\sigma,\gamma)$ and state summaries.
\begin{enumerate}[label=(\arabic*)]
\item Initialize $\vect{\theta}_{0:T}$, $\vect{v}$, $\vect{s}$, $\sigma$, and
$\gamma$.
\item For $m=1,\ldots,N_{\mathrm{burn}}+N_{\mathrm{mcmc}}$:
  \begin{enumerate}[label=(\alph*)]
  \item sample $v_t$ from \eqref{eq:dyn_exdqlm_v_cond}, $t=1,\ldots,T$;
  \item sample $s_t$ from \eqref{eq:dyn_exdqlm_s_cond}, $t=1,\ldots,T$;
  \item sample $\vect{\theta}_{0:T}$ by FFBS using
  Section~\ref{sec:gaussian_state_update} and the exDQLM MCMC row of
  Table~\ref{tab:gaussian_state_update_map};
  \item sample $\sigma$ from \eqref{eq:dyn_exdqlm_sigma_cond};
  \item sample $\gamma$ on $(L,U)$ with the bounded slice sampler in
  Section~\ref{sec:slice}, targeting \eqref{eq:dyn_exdqlm_gamma_kernel};
  \item if $m>N_{\mathrm{burn}}$, retain the current draw.
  \end{enumerate}
\item Return the retained draws and state summaries.
\end{enumerate}
\end{suppalgorithm}
Setting $\gamma=0$ drops the $s_t$ block and reduces the sampler to the
AL/DQLM kernel in Algorithm~\ref{alg:dyn-dqlm-mcmc}.

\subsection{Laplace--delta VB approximation}

The mean-field family is
\begin{equation}
q(\vect{\theta}_{0:T},\vect{v},\vect{s},\sigma,\gamma)
=q(\vect{\theta}_{0:T})\prod_{t=1}^Tq(v_t)\prod_{t=1}^Tq(s_t)\,q(\sigma,\gamma).
\label{eq:dyn_exdqlm_mf}
\end{equation}
The conjugate factors are updated conditionally on moments from
$q(\sigma,\gamma)$. Define
\begin{align}
\kappa_1 &= \E\{1/(\sigma B_\gamma)\},&
\kappa_2 &= \E\{A_\gamma/(\sigma B_\gamma)\},&
\kappa_3 &= \E\{A_\gamma^2/(\sigma B_\gamma)\},\\
\kappa_4 &= \E\{C_\gamma|\gamma|/B_\gamma\},&
\kappa_5 &= \E\{\sigma C_\gamma^2\gamma^2/B_\gamma\},&
\kappa_6 &= \E\{A_\gamma C_\gamma|\gamma|/B_\gamma\}.
\label{eq:dyn_exdqlm_kappa}
\end{align}
Let $m_{1/v,t}=\E(v_t^{-1})$, $\bar s_t=\E(s_t)$, and
$\overline{s_t^2}=\E(s_t^2)$. The state factor is Gaussian and is updated by
the deterministic smoother in Section~\ref{sec:gaussian_state_update} with
observation precision
\begin{equation}
\phi_t=\kappa_1m_{1/v,t}
\label{eq:dyn_exdqlm_vb_precision}
\end{equation}
and pseudo-response
\begin{equation}
y_t^{\mathrm{VB}}
=
\frac{
y_t\phi_t-\kappa_2-\kappa_4\bar s_t m_{1/v,t}
}{\phi_t}.
\label{eq:dyn_exdqlm_vb_pseudo}
\end{equation}
Equivalently, in the original-data form
\eqref{eq:generic_gaussian_obs}, the offset is
$d_t=\{\kappa_2+\kappa_4\bar s_t m_{1/v,t}\}/\phi_t$ and the variance is
$Q_t=\phi_t^{-1}$. To see this, let
$\eta_t=\mat{F}_t^\top\vect{\theta}_t$ and let $\E_{-\theta}$ denote
expectation over all variational factors except the state factor. The terms
in the expected quadratic loss that depend on $\eta_t$ are
\begin{align}
&\E_{-\theta}\left[
\frac{
\{y_t-\eta_t-A_\gamma v_t-C_\gamma\sigma|\gamma|s_t\}^2
}{
\sigma B_\gamma v_t
}
\right] \nonumber\\
&\qquad =
\phi_t(y_t-\eta_t)^2
-
2\{\kappa_2+\kappa_4\bar s_t m_{1/v,t}\}(y_t-\eta_t)
+\text{constant},
\label{eq:dyn_exdqlm_ldvb_state_square}
\end{align}
so completing the square gives \eqref{eq:dyn_exdqlm_vb_pseudo} and
\eqref{eq:dyn_exdqlm_vb_precision}.
If $\eta_t=\mat{F}_t^\top\vect{\theta}_t$, write
$m_{\eta,t}=\E(\eta_t)$, $S_{\eta,t}=\Var(\eta_t)$,
$m_{\Delta,t}=y_t-m_{\eta,t}$, and
$S_{\Delta,t}=m_{\Delta,t}^2+S_{\eta,t}$. Then
\begin{align}
q(v_t)&=\GIG\left(
\frac{1}{2},\
\kappa_1S_{\Delta,t}-2\kappa_4m_{\Delta,t}\bar s_t+\kappa_5\overline{s_t^2},\
\kappa_3+2\E(\sigma^{-1})
\right),\\
q(s_t)&=\TNp(\mu_{s,t}^{\mathrm{VB}},V_{s,t}^{\mathrm{VB}}),\\
V_{s,t}^{\mathrm{VB}}
&=\left(1+\kappa_5m_{1/v,t}\right)^{-1},\\
\mu_{s,t}^{\mathrm{VB}}
&=V_{s,t}^{\mathrm{VB}}
\left(\kappa_4m_{\Delta,t}m_{1/v,t}-\kappa_6\right).
\label{eq:dyn_exdqlm_vb_vs}
\end{align}
The nonconjugate factor $q(\sigma,\gamma)$ is approximated by a
Laplace--delta step on transformed coordinates
\begin{equation}
\ell_\sigma=\log\sigma,\qquad
z_\gamma=\logit\left(\frac{\gamma-L}{U-L}\right).
\label{eq:dyn_exdqlm_ld_transform}
\end{equation}
The transformed objective is the expected complete log-joint plus the
Jacobian. Here $\E_{-(\sigma,\gamma)}$ denotes expectation with respect to all
variational factors except $q(\sigma,\gamma)$:
\begin{align}
h_z(\ell_\sigma,z_\gamma)
&=
\E_{-(\sigma,\gamma)}\{\log p(\vect{y},\vect{\theta}_{0:T},
\vect{v},\vect{s},\sigma,\gamma)\} \nonumber\\
&\quad
+\ell_\sigma+\log(U-L)+\log r+\log(1-r),
\qquad r=\expit(z_\gamma).
\label{eq:dyn_exdqlm_ld_objective}
\end{align}
The fitted object returns the transformed mode, the approximate covariance
matrix, delta-method moments, and ELBO components.

\begin{suppalgorithm}{Dynamic exDQLM LDVB}
\label{alg:dyn-exdqlm-ldvb}
\par\noindent\emph{Input:} observations $\vect{y}$, model matrices
$(\mat{F}_t,\mat{G}_t,\mat{W}_t)$, initial state prior, exAL prior
hyperparameters, admissible interval $(L,U)$ for $\gamma$, convergence
tolerance, and optional posterior-sampling control \code{n.samp}.
\par\noindent\emph{Output:} fitted variational factors for
$(\vect{\theta}_{0:T},\vect{v},\vect{s},\sigma,\gamma)$, convergence
diagnostics, ELBO values when requested, and optional approximate posterior
draws.
\begin{enumerate}[label=(\arabic*)]
\item Initialize variational parameters for $q(\vect{\theta}_{0:T})$,
$q(v_t)$, $q(s_t)$, and $q(\sigma,\gamma)$.
\item Repeat until the convergence criterion is met:
  \begin{enumerate}[label=(\alph*)]
  \item update $q(\vect{\theta}_{0:T})$ by the deterministic smoother in
  Section~\ref{sec:gaussian_state_update}, using the exDQLM LDVB row of
  Table~\ref{tab:gaussian_state_update_map};
  \item update $q(v_t)$ and $q(s_t)$ using \eqref{eq:dyn_exdqlm_vb_vs};
  \item update $q(\sigma,\gamma)$ on the transformed coordinates
  $(\ell_\sigma,z_\gamma)$ in \eqref{eq:dyn_exdqlm_ld_transform} by the
  Laplace--delta approximation to \eqref{eq:dyn_exdqlm_ld_objective};
  \item monitor moment changes and, when ELBO monitoring is enabled, $\elbo(q)$.
  \end{enumerate}
\item Draw approximate posterior samples from the fitted variational factors
when requested by \code{n.samp}.
\item Return the fitted factors, diagnostics, and optional draws.
\end{enumerate}
\end{suppalgorithm}
In the AL special case, the nonconjugate $(\sigma,\gamma)$ block is replaced
by the inverse-gamma $q(\sigma)$ update in Algorithm~\ref{alg:dyn-dqlm-vb}.
This exAL path is the default in \code{exdqlmMCMC()} and \code{exdqlmLDVB()}
when \code{dqlm.ind = FALSE}.

\paragraph{ELBO.}
The LDVB ELBO has the block decomposition
\begin{equation}
\elbo(q)=
\elbo_{\mathrm{like}}+\elbo_\theta+\elbo_v+\elbo_s+\elbo_{\sigma\gamma}
+\sum_{t=1}^TH\{q(v_t)\}
+\sum_{t=1}^TH\{q(s_t)\}
+H\{q(\sigma,\gamma)\}.
\label{eq:dyn_exdqlm_elbo}
\end{equation}
The state block is
\begin{equation}
\elbo_\theta
=\E_q\left\{\log p(\vect{\theta}_0)
+\sum_{t=1}^T\log p(\vect{\theta}_t\mid\vect{\theta}_{t-1})\right\}
+H\{q(\vect{\theta}_{0:T})\}.
\label{eq:dyn_exdqlm_elbo_theta}
\end{equation}
As in the AL case, this Gaussian state-space contribution can be evaluated by
the auxiliary Kalman normalizing-constant identity for the current
pseudo-response and variance.

Let
$\bar\ell_\sigma=\E(\log\sigma)$,
$\bar\ell_B=\E(\log B_\gamma)$,
$\bar\ell_{\pi}=\E\{\log\pigamma(\gamma)\}$,
$m_{v,t}=\E(v_t)$, $m_{1/v,t}=\E(v_t^{-1})$,
$m_{\log v,t}=\E(\log v_t)$, $\bar s_t=\E(s_t)$, and
$\overline{s_t^2}=\E(s_t^2)$. Define
\begin{align}
\mathcal{Q}_t
&=
\kappa_1S_{\Delta,t}m_{1/v,t}
-2\kappa_2m_{\Delta,t}
+\kappa_3m_{v,t}
-2\kappa_4m_{\Delta,t}\bar s_t m_{1/v,t} \nonumber\\
&\quad
+\kappa_5\overline{s_t^2}m_{1/v,t}
+2\kappa_6\bar s_t.
\label{eq:dyn_exdqlm_elbo_mathcalQ}
\end{align}
Then the likelihood, latent-prior, and nonconjugate-prior blocks are
\begin{align}
\elbo_{\mathrm{like}}
&=
-\frac{T}{2}\log(2\pi)
-\frac{T}{2}(\bar\ell_\sigma+\bar\ell_B)
-\frac{1}{2}\sum_{t=1}^Tm_{\log v,t}
-\frac{1}{2}\sum_{t=1}^T\mathcal{Q}_t,\\
\elbo_v
&=
-T\bar\ell_\sigma-\E(\sigma^{-1})\sum_{t=1}^Tm_{v,t},\\
\elbo_s
&=
T\log 2-\frac{T}{2}\log(2\pi)
-\frac{1}{2}\sum_{t=1}^T\overline{s_t^2},\\
\elbo_{\sigma\gamma}
&=
a_\sigma\log b_\sigma-\log\Gamma(a_\sigma)
-(a_\sigma+1)\bar\ell_\sigma
-b_\sigma\E(\sigma^{-1})
+\bar\ell_{\pi}.
\label{eq:dyn_exdqlm_elbo_terms}
\end{align}
The GIG and truncated-normal entropy terms are given by
\eqref{eq:entropy_gig} and \eqref{eq:entropy_tnp}. The
$(\sigma,\gamma)$ entropy and the smooth expectations in the $\kappa_j$,
$\bar\ell_B$, $\bar\ell_\pi$, and $\E(\sigma^{-1})$ terms are evaluated under
the Laplace--delta approximation in \eqref{eq:entropy_ld} and
\eqref{eq:delta_generic}.

\section{Transfer-function exDQLM by state augmentation}
\label{sec:transfer_function}

Transfer-function inference uses the same AL/exAL likelihood and latent-variable
updates as the dynamic models above. The only change is the state-space
construction. Let $x_t\in\R^k$ be the input vector, $\zeta_t$ the transfer
state, and $\psi_t\in\R^k$ the input-effect state. For fixed transfer rate
$\lambda$, the augmented model uses
\begin{align}
\widetilde{\mat{F}}_t^\top
&=(\mat{F}_t^\top,\ 1,\ 0_k^\top),&
\widetilde{\vect{\theta}}_t^\top
&=(\vect{\theta}_t^\top,\ \zeta_t,\ \psi_t^\top),
\label{eq:tf_aug_state}\\
\widetilde{\mat{G}}_t
&=
\operatorname{blockdiag}\left\{
\mat{G}_t,\,
\begin{pmatrix}
\lambda & x_t^\top\\
0_k & I_k
\end{pmatrix}
\right\},&
\widetilde{\mat{W}}_t
&=\operatorname{blockdiag}\{\mat{W}_t^\theta,\omega_t,\mat{W}_t^\psi\}.
\label{eq:tf_aug_mats}
\end{align}
The prior for $(\zeta_0,\psi_0^\top)^\top$ is
$\N(m_0^{\mathrm{tf}},C_0^{\mathrm{tf}})$ and is appended to the base
state prior. Conditional on $\lambda$, the augmented model has the same
posterior form as the base exDQLM with state
$\widetilde{\vect{\theta}}_t$. Consequently,
the Gaussian state update in Section~\ref{sec:gaussian_state_update}, together
with Algorithms~\ref{alg:dyn-exdqlm-mcmc} and~\ref{alg:dyn-exdqlm-ldvb},
applies after replacing
$(\mat{F}_t,\mat{G}_t,\mat{W}_t,\vect{\theta}_t)$ by
$(\widetilde{\mat{F}}_t,\widetilde{\mat{G}}_t,\widetilde{\mat{W}}_t,
\widetilde{\vect{\theta}}_t)$.
The current transfer-fitting routines treat $\lambda$ as fixed by the user or
by an external grid search; they do not sample or optimize $\lambda$
internally.
The package inputs \code{X}, \code{lam}, \code{tf.df}, \code{tf.m0}, and
\code{tf.C0} correspond to $x_t$, $\lambda$, the transfer discount factors,
$m_0^{\mathrm{tf}}$, and $C_0^{\mathrm{tf}}$; \code{median.kt} is returned as a
summary of the fitted transfer response.

\begin{suppalgorithm}{Transfer-function state augmentation and fitting}
\label{alg:transfer-fitting}
\par\noindent\emph{Input:} base dynamic model, transfer covariates $x_t$,
fixed transfer rate $\lambda$, transfer discount factors, and transfer-state
prior $(m_0^{\mathrm{tf}},C_0^{\mathrm{tf}})$.
\par\noindent\emph{Output:} fitted base-state and transfer-state summaries,
including the transfer response summary \code{median.kt}.
\begin{enumerate}[label=(\arabic*)]
\item Construct $\widetilde{\mat{F}}_t$, $\widetilde{\mat{G}}_t$,
$\widetilde{\mat{W}}_t$, and $\widetilde{\vect{\theta}}_t$ from
\eqref{eq:tf_aug_state}--\eqref{eq:tf_aug_mats}.
\item Set the augmented initial prior.
\item Run the selected fitting routine on the augmented state-space system:
MCMC uses Algorithm~\ref{alg:dyn-exdqlm-mcmc} and LDVB uses
Algorithm~\ref{alg:dyn-exdqlm-ldvb}.
\item Return summaries for both the base state and the transfer states.
\end{enumerate}
\end{suppalgorithm}
The transfer-function extension changes only the state-space matrices; the
observation likelihood and latent-variable updates are unchanged.

\section{Static AL and exAL regression with Gaussian priors}
\label{sec:static_ridge}

\subsection{Model hierarchy and posterior target}

Let $\mat{X}\in\R^{n\times p}$ have row vectors $\vect{x}_i^\top$, and let
$\vect{\beta}\in\R^p$. The ridge/Gaussian coefficient prior is
\begin{equation}
\vect{\beta}\sim\N(\vect{b}_{\beta,0},\mat{V}_{\beta,0}),\qquad
\vect{b}_{\beta,0}\in\R^p,\quad \mat{V}_{\beta,0}\in\R^{p\times p}.
\label{eq:static_beta_ridge}
\end{equation}
In the package interface, the same Gaussian update is used for
\code{beta_prior = "ridge"}. When \code{beta_prior = "rhs\_ns"}, the diagonal
prior precision in this update is replaced by the \code{rhs\_ns} expected precision
described in Section~\ref{sec:static_rhs}. Thus this supplement covers the
Gaussian/ridge prior and the conjugate \code{rhs\_ns} prior, both for MCMC and LDVB.
Under the exAL likelihood,
\begin{align}
v_i\mid\sigma &\sim \Exp(\text{rate}=1/\sigma),\qquad
s_i\sim\TNp(0,1),\\
y_i\mid\vect{\beta},\sigma,\gamma,v_i,s_i
&\sim
\N\left(
\vect{x}_i^\top\vect{\beta}+C_\gamma\sigma|\gamma|s_i+A_\gamma v_i,\
\sigma B_\gamma v_i
\right).
\label{eq:static_exal_hierarchy}
\end{align}
The static AL model is obtained by setting $\gamma=0$, removing $s_i$, and
using $A_{p_0},B_{p_0}$. The exAL posterior target is
\begin{align}
p(\vect{\beta},\sigma,\gamma,\vect{v},\vect{s}\mid\vect{y})
&\propto
p(\vect{\beta})p(\sigma)\pigamma(\gamma)
\prod_{i=1}^n
p(y_i\mid\vect{\beta},\sigma,\gamma,v_i,s_i)
p(v_i\mid\sigma)p(s_i).
\label{eq:static_exal_target}
\end{align}
Unlike the dynamic models in Sections~\ref{sec:dyn_dqlm}--\ref{sec:transfer_function},
the static regression update does not use Kalman filtering, FFBS, or RTS
smoothing. The coefficient block is a Gaussian weighted-regression update.

\subsection{MCMC updates}

For fixed $(\sigma,\gamma,\vect{v},\vect{s})$, define
\begin{equation}
y_i^\star=y_i-C_\gamma\sigma|\gamma|s_i-A_\gamma v_i,\qquad
w_i=(\sigma B_\gamma v_i)^{-1}.
\end{equation}
Let $\mat{W}=\diag(w_1,\ldots,w_n)$ and
$\vect{y}^\star=(y_1^\star,\ldots,y_n^\star)^\top$. Then
\begin{align}
\mat{\Sigma}_{\beta\mid\cdot}
&=\left(\mat{V}_{\beta,0}^{-1}+\mat{X}^\top\mat{W}\mat{X}\right)^{-1},\\
\vect{m}_{\beta\mid\cdot}
&=\mat{\Sigma}_{\beta\mid\cdot}
\left(\mat{V}_{\beta,0}^{-1}\vect{b}_{\beta,0}+\mat{X}^\top\mat{W}\vect{y}^\star\right),\\
\vect{\beta}\mid\cdot
&\sim\N(\vect{m}_{\beta\mid\cdot},\mat{\Sigma}_{\beta\mid\cdot}).
\label{eq:static_beta_cond}
\end{align}
Let $D_\gamma=C_\gamma|\gamma|$. The remaining exAL full conditionals are
\begin{align}
s_i\mid\cdot
&\sim \TNp(\mu_{s,i},V_{s,i}),\\
V_{s,i}
&=\left(1+\frac{\sigma D_\gamma^2}{B_\gamma v_i}\right)^{-1},\\
\mu_{s,i}
&=V_{s,i}\frac{D_\gamma(y_i-\vect{x}_i^\top\vect{\beta}-A_\gamma v_i)}
{B_\gamma v_i},
\label{eq:static_exal_s_cond}
\end{align}
and, with $r_i=y_i-\vect{x}_i^\top\vect{\beta}-\sigma D_\gamma s_i$,
\begin{equation}
v_i\mid\cdot
\sim
\GIG\left(
\frac{1}{2},\
\frac{r_i^2}{\sigma B_\gamma},\
\frac{A_\gamma^2}{\sigma B_\gamma}+\frac{2}{\sigma}
\right).
\label{eq:static_exal_v_cond}
\end{equation}
For fixed $\gamma$, define $t_i=y_i-\vect{x}_i^\top\vect{\beta}-A_\gamma v_i$.
Then
\begin{equation}
\sigma\mid\cdot
\sim
\GIG\left(
-a_\sigma-\frac{3n}{2},\
2b_\sigma+2\sum_{i=1}^n v_i+\sum_{i=1}^n\frac{t_i^2}{B_\gamma v_i},\
\sum_{i=1}^n\frac{D_\gamma^2s_i^2}{B_\gamma v_i}
\right).
\label{eq:static_exal_sigma_cond}
\end{equation}
The nonconjugate conditional kernel for $\gamma$ is
\begin{align}
\pi(\gamma\mid\cdot)
&\propto
\pigamma(\gamma)\,[B_\gamma]^{-n/2}
\exp\left\{
\sum_{i=1}^n\frac{D_\gamma s_it_i}{B_\gamma v_i}
-\frac{1}{2\sigma}\sum_{i=1}^n\frac{t_i^2}{B_\gamma v_i}
-\frac{\sigma}{2}\sum_{i=1}^n\frac{D_\gamma^2s_i^2}{B_\gamma v_i}
\right\},
\label{eq:static_exal_gamma_cond}
\end{align}
on $(L,U)$. The package default updates this one-dimensional target by the
bounded slice sampler in Section~\ref{sec:slice}.

For the AL special case, the Gibbs updates reduce to
\begin{align}
\vect{\beta}\mid\cdot
&\sim\N(\vect{m}_{\beta\mid\cdot},\mat{\Sigma}_{\beta\mid\cdot}),
\label{eq:static_al_beta_cond}\\
\sigma\mid\cdot
&\sim
\IG\left(
a_\sigma+\frac{3n}{2},\
b_\sigma+\sum_{i=1}^n v_i+
\sum_{i=1}^n\frac{(y_i-\vect{x}_i^\top\vect{\beta}-A_{p_0}v_i)^2}{2B_{p_0}v_i}
\right),
\label{eq:static_al_sigma_cond}\\
v_i\mid\cdot
&\sim
\GIG\left(
\frac{1}{2},\
\frac{(y_i-\vect{x}_i^\top\vect{\beta})^2}{\sigma B_{p_0}},\
\frac{A_{p_0}^2}{\sigma B_{p_0}}+\frac{2}{\sigma}
\right).
\label{eq:static_al_v_cond}
\end{align}

\begin{suppalgorithm}{Static AL/exAL MCMC}
\label{alg:static-mcmc}
\par\noindent\emph{Input:} response vector $\vect{y}$, design matrix
$\mat{X}$, target quantile $p_0$, coefficient prior specification, AL/exAL
prior hyperparameters, and MCMC controls $N_{\mathrm{burn}}$ and
$N_{\mathrm{mcmc}}$.
\par\noindent\emph{Output:} retained posterior draws for regression
coefficients, latent variables, scale parameters, and shrinkage parameters when
\code{beta_prior = "rhs\_ns"}.
\begin{enumerate}[label=(\arabic*)]
\item Initialize $\vect{\beta}$, $\vect{v}$, $\vect{s}$, $\sigma$, and
$\gamma$; omit $\vect{s}$ and fix $\gamma=0$ for the AL path.
\item For $m=1,\ldots,N_{\mathrm{burn}}+N_{\mathrm{mcmc}}$:
  \begin{enumerate}[label=(\alph*)]
  \item sample $\vect{\beta}$ from \eqref{eq:static_beta_cond};
  \item sample $v_i$ from \eqref{eq:static_exal_v_cond},
  $i=1,\ldots,n$, with the AL simplification in \eqref{eq:static_al_v_cond};
  \item for exAL, sample $s_i$ from \eqref{eq:static_exal_s_cond},
  $i=1,\ldots,n$;
  \item sample $\sigma$ from \eqref{eq:static_exal_sigma_cond}, with the AL
  simplification in \eqref{eq:static_al_sigma_cond};
  \item for exAL, sample $\gamma$ on $(L,U)$ with the bounded slice sampler,
  targeting \eqref{eq:static_exal_gamma_cond};
  \item if \code{beta_prior = "rhs\_ns"}, update the \code{rhs\_ns} scale block in
  \eqref{eq:rhs_ns_mcmc};
  \item if $m>N_{\mathrm{burn}}$, retain the current draw.
  \end{enumerate}
\item Return the retained posterior draws.
\end{enumerate}
\end{suppalgorithm}

\subsection{Laplace--delta VB updates and ELBO}

The static exAL variational family is
\begin{equation}
q(\vect{\beta},\vect{v},\vect{s},\sigma,\gamma)
=q(\vect{\beta})\prod_{i=1}^nq(v_i)\prod_{i=1}^nq(s_i)q(\sigma,\gamma).
\label{eq:static_exal_mf}
\end{equation}
The coefficient factor is Gaussian. With the same moment definitions as in
\eqref{eq:dyn_exdqlm_kappa},
\begin{align}
\omega_i&=\kappa_1\E(v_i^{-1}),\\
\eta_i&=
y_i\kappa_1\E(v_i^{-1})
-\kappa_4\E(v_i^{-1})\E(s_i)
-\kappa_2,\\
\mat{\Sigma}_\beta
&=\left(\mat{V}_{\beta,0}^{-1}+\mat{X}^\top\diag(\omega_1,\ldots,\omega_n)\mat{X}\right)^{-1},\\
\vect{m}_\beta
&=\mat{\Sigma}_\beta
\left(\mat{V}_{\beta,0}^{-1}\vect{b}_{\beta,0}+\mat{X}^\top\vect{\eta}\right),
\label{eq:static_exal_vb_beta}
\end{align}
where $\vect{\eta}=(\eta_1,\ldots,\eta_n)^\top$. The $q(v_i)$ and $q(s_i)$
updates are also closed form conditional on moments of $q(\sigma,\gamma)$.
Let
\begin{equation}
\bar r_i=y_i-\vect{x}_i^\top\vect{m}_\beta,\qquad
\overline{r_i^2}=\bar r_i^2+\vect{x}_i^\top\mat{\Sigma}_\beta\vect{x}_i.
\end{equation}
Then
\begin{align}
q(v_i)
&=
\GIG\left(
\frac{1}{2},\
\kappa_1\overline{r_i^2}
-2\kappa_4\bar r_i\E(s_i)
+\kappa_5\E(s_i^2),\
\kappa_3+2\E(\sigma^{-1})
\right),\\
q(s_i)&=\TNp(\mu_{s,i}^{\mathrm{VB}},V_{s,i}^{\mathrm{VB}}),\\
V_{s,i}^{\mathrm{VB}}
&=\left(1+\kappa_5\E(v_i^{-1})\right)^{-1},\\
\mu_{s,i}^{\mathrm{VB}}
&=V_{s,i}^{\mathrm{VB}}
\left\{\kappa_4\bar r_i\E(v_i^{-1})-\kappa_6\right\}.
\label{eq:static_exal_vb_vs}
\end{align}
The nonconjugate factor $q(\sigma,\gamma)$ is proportional to
\begin{align}
q(\sigma,\gamma)
&\propto
\pigamma(\gamma)[B_\gamma]^{-n/2}
\sigma^{-(a_\sigma+1+3n/2)}
\exp\left\{
-\frac{1}{2}\left(\frac{\bar c(\gamma)}{\sigma}
+\bar d(\gamma)\sigma\right)+\bar e(\gamma)
\right\},
\label{eq:static_exal_q_sigmagam}
\end{align}
where
\begin{align}
\bar c(\gamma)
&=2b_\sigma+2\sum_{i=1}^n\E(v_i)
+\sum_{i=1}^n
\frac{\E\{(y_i-\vect{x}_i^\top\vect{\beta}-A_\gamma v_i)^2/v_i\}}
{B_\gamma},\\
\bar d(\gamma)
&=\sum_{i=1}^n\frac{C_\gamma^2\gamma^2}{B_\gamma}\E\{s_i^2/v_i\},\\
\bar e(\gamma)
&=\sum_{i=1}^n
\frac{C_\gamma|\gamma|}{B_\gamma}
\E\{s_i(y_i-\vect{x}_i^\top\vect{\beta}-A_\gamma v_i)/v_i\}.
\end{align}
This factor is approximated by the same Laplace--delta method as in the
dynamic exDQLM. The AL path sets $\gamma=0$ and uses
$q(\vect{\beta})q(\sigma)\prod_iq(v_i)$ with the same closed-form updates as
Section~\ref{sec:dyn_dqlm}, replacing state smoothing moments by
$\E\{(y_i-\vect{x}_i^\top\vect{\beta})^2\}$.

\begin{suppalgorithm}{Static AL/exAL LDVB}
\label{alg:static-ldvb}
\par\noindent\emph{Input:} response vector $\vect{y}$, design matrix
$\mat{X}$, target quantile $p_0$, coefficient prior specification, AL/exAL
prior hyperparameters, convergence tolerance, and optional posterior-sampling
control \code{n.samp}.
\par\noindent\emph{Output:} fitted variational factors, convergence
diagnostics, ELBO values, and optional approximate posterior draws.
\begin{enumerate}[label=(\arabic*)]
\item Initialize $q(\vect{\beta})$, the latent-variable factors, the scale
factors, and the $(\sigma,\gamma)$ block.
\item Repeat until the convergence criterion is met:
  \begin{enumerate}[label=(\alph*)]
  \item update $q(\vect{\beta})$ using \eqref{eq:static_exal_vb_beta};
  \item update $q(v_i)$ and, for exAL, $q(s_i)$ using
  \eqref{eq:static_exal_vb_vs};
  \item update $q(\sigma)$ for AL or $q(\sigma,\gamma)$ for exAL using the
  Laplace--delta approximation to \eqref{eq:static_exal_q_sigmagam};
  \item if \code{beta_prior = "rhs\_ns"}, update the \code{rhs\_ns} factors in
  \eqref{eq:rhs_ns_vb_updates};
  \item monitor moment changes and $\elbo(q)$.
  \end{enumerate}
\item Draw approximate posterior samples when requested by \code{n.samp}.
\item Return the fitted factors, diagnostics, and optional draws.
\end{enumerate}
\end{suppalgorithm}
These static updates are implemented in \code{exalStaticMCMC()} and
\code{exalStaticLDVB()}, with \code{beta_prior} selecting the Gaussian/ridge
or \code{rhs\_ns} coefficient prior.

The static exAL ELBO is computed from the same blocks as the dynamic exDQLM,
with the Gaussian state factor replaced by the Gaussian coefficient factor.
Let
$\bar\ell_\sigma=\E(\log\sigma)$,
$\bar\ell_B=\E(\log B_\gamma)$,
$\bar\ell_\pi=\E\{\log\pigamma(\gamma)\}$,
$m_{v,i}=\E(v_i)$, $m_{1/v,i}=\E(v_i^{-1})$,
$m_{\log v,i}=\E(\log v_i)$, $\bar s_i=\E(s_i)$, and
$\overline{s_i^2}=\E(s_i^2)$. Also define
\begin{equation}
t_i=y_i-\vect{x}_i^\top\vect{m}_\beta,\qquad
q_i=\vect{x}_i^\top\mat{\Sigma}_\beta\vect{x}_i.
\end{equation}
Then
\begin{align}
\elbo(q)
&=
\elbo_{\mathrm{like}}+\elbo_v+\elbo_s+\elbo_\beta
+\elbo_{\sigma\gamma}
+H\{q(\vect{\beta})\}
+\sum_{i=1}^nH\{q(v_i)\}\nonumber\\
&\quad
+\sum_{i=1}^nH\{q(s_i)\}
+H\{q(\sigma,\gamma)\},
\label{eq:static_exal_elbo}
\end{align}
where
\begin{align}
\elbo_{\mathrm{like}}
&=
-\frac{n}{2}\log(2\pi)
-\frac{n}{2}(\bar\ell_B+\bar\ell_\sigma)
-\frac{1}{2}\sum_{i=1}^nm_{\log v,i}\nonumber\\
&\quad
-\frac{1}{2}\sum_{i=1}^n
\left\{
\kappa_1m_{1/v,i}(t_i^2+q_i)
-2\kappa_2t_i+\kappa_3m_{v,i}
\right\}\nonumber\\
&\quad
+\sum_{i=1}^n
\left\{
\kappa_4\bar s_i m_{1/v,i}t_i
-\kappa_6\bar s_i
-\frac{1}{2}\kappa_5\overline{s_i^2}m_{1/v,i}
\right\},\\
\elbo_v
&=
-n\bar\ell_\sigma-\E(\sigma^{-1})\sum_{i=1}^nm_{v,i},\\
\elbo_s
&=
n\log 2-\frac{n}{2}\log(2\pi)
-\frac{1}{2}\sum_{i=1}^n\overline{s_i^2},\\
\elbo_{\sigma\gamma}
&=
a_\sigma\log b_\sigma-\log\Gamma(a_\sigma)
-(a_\sigma+1)\bar\ell_\sigma
-b_\sigma\E(\sigma^{-1})
+\bar\ell_\pi,\\
\elbo_\beta
&=\E_q\{\log p(\vect{\beta})\}.
\label{eq:static_exal_elbo_terms}
\end{align}
For the ridge prior, $\elbo_\beta$ is the Gaussian prior expectation. For the
\code{rhs\_ns} prior, it is replaced by the \code{rhs\_ns} contribution in
Section~\ref{sec:static_rhs}. The Gaussian, GIG, truncated-normal, and
Laplace--delta entropy terms are given by
\eqref{eq:entropy_gaussian}, \eqref{eq:entropy_gig},
\eqref{eq:entropy_tnp}, and \eqref{eq:entropy_ld}, respectively. The AL path
is recovered by setting $\gamma=0$ and dropping the $s_i$ and
$(\sigma,\gamma)$ Laplace--delta blocks; then $q(\sigma)$ is inverse-gamma and
uses \eqref{eq:entropy_ig}.

\section{Static regression with the Nishimura--Suchard 
regularized horseshoe prior}
\label{sec:static_rhs}

This section focuses on the \code{rhs\_ns} coefficient prior, the
conditionally conjugate regularized horseshoe variant used for sparse static
regression. Users only need this section when fitting static models with
\code{beta_prior = "rhs\_ns"}. The prior enters the static Gaussian update for
$\vect{\beta}$ through a prior precision matrix. Let
$\mathcal{A}\subseteq\{1,\ldots,p\}$ denote the active shrinkage set; by
default an intercept, if present, is not shrunk. For $j\notin\mathcal{A}$ the
package uses a fixed intercept precision. The direct \code{rhs} path is not
developed here because its log-scale block is nonconjugate; the
\code{rhs\_ns} path is the conjugate sparse-prior formulation documented in
this supplement.

\subsection{Nishimura--Suchard prior hierarchy}

For $j\in\mathcal{A}$, the \code{rhs\_ns} hierarchy is
\begin{align}
\beta_j\mid\tau^2,\lambda_j^2,\zeta^2
&\propto
\N(\beta_j\mid0,\tau^2\lambda_j^2)
\N(\beta_j\mid0,\zeta^2),\\
\lambda_j^2\mid\nu_j &\sim \IG(1/2,1/\nu_j),&
\nu_j&\sim\IG(1/2,1),\\
\tau^2\mid\xi &\sim\IG(1/2,1/\xi),&
\xi&\sim\IG(1/2,1/\tau_0^2),\\
\zeta^2&\sim\IG(a_\zeta,b_\zeta),
\label{eq:rhs_ns_hierarchy}
\end{align}
unless $\zeta^2$ is fixed by the user. The implementation includes the
scale-dependent Gaussian log-normalizing terms shown in the ELBO contribution
below. Thus the proportionality in \eqref{eq:rhs_ns_hierarchy} denotes the
coefficient kernel, while the surrounding scale updates and ELBO use the
corresponding normal-density contributions. The resulting prior precision
contribution for active coefficients is
\begin{equation}
\omega^{\mathrm{prior}}_j=\frac{1}{\tau^2\lambda_j^2}+\frac{1}{\zeta^2}.
\label{eq:rhs_ns_precision}
\end{equation}
Consequently, the $\vect{\beta}$ update in MCMC or LDVB is obtained from the
static Gaussian update after replacing $\mat{V}_{\beta,0}^{-1}$ by a diagonal precision
matrix with active entries $\omega^{\mathrm{prior}}_j$ or their variational
expectations.

\subsection{Nishimura--Suchard MCMC updates}

Conditioning on $\vect{\beta}$, the shrinkage block has closed-form
inverse-gamma updates:
\begin{align}
\lambda_j^2\mid\cdot
&\sim
\IG\left(1,\ \frac{1}{\nu_j}+\frac{\beta_j^2}{2\tau^2}\right),\\
\nu_j\mid\cdot
&\sim
\IG\left(1,\ 1+\frac{1}{\lambda_j^2}\right),\\
\tau^2\mid\cdot
&\sim
\IG\left(\frac{|\mathcal{A}|+1}{2},\
\frac{1}{\xi}+\frac{1}{2}\sum_{j\in\mathcal{A}}\frac{\beta_j^2}{\lambda_j^2}
\right),\\
\xi\mid\cdot
&\sim
\IG\left(1,\ \frac{1}{\tau_0^2}+\frac{1}{\tau^2}\right),\\
\zeta^2\mid\cdot
&\sim
\IG\left(a_\zeta+\frac{|\mathcal{A}|}{2},\
b_\zeta+\frac{1}{2}\sum_{j\in\mathcal{A}}\beta_j^2
\right),
\label{eq:rhs_ns_mcmc}
\end{align}
with the last line omitted when $\zeta^2$ is fixed.
These equations are the conditional targets for the \code{rhs\_ns} scale block. The
implementation may temporarily freeze or schedule the global-scale update during
configured warmup iterations for numerical stability; this changes only the
update schedule, not the conditional target displayed above.

\subsection{Nishimura--Suchard variational updates and ELBO contribution}

The package uses a mean-field approximation over the \code{rhs\_ns} scale block. For
$j\in\mathcal{A}$,
\begin{align}
q(\lambda_j^2)&=\IG(a_{\lambda j},b_{\lambda j}),&
q(\nu_j)&=\IG(a_{\nu j},b_{\nu j}),\\
q(\tau^2)&=\IG(a_\tau,b_\tau),&
q(\xi)&=\IG(a_\xi,b_\xi),\\
q(\zeta^2)&=\IG(a_\zeta^\star,b_\zeta^\star),
\end{align}
with $a_{\lambda j}=a_{\nu j}=a_\xi=1$ and
$a_\tau=(|\mathcal{A}|+1)/2$. Let
$\bar\beta_j^2=\E_q(\beta_j^2)$. The coordinate updates are
\begin{align}
b_{\lambda j}
&=\frac{1}{2}\bar\beta_j^2\,\E\{(\tau^2)^{-1}\}+\E(\nu_j^{-1}),\\
b_{\nu j}
&=1+\E\{(\lambda_j^2)^{-1}\},\\
b_\tau
&=\frac{1}{2}\sum_{j\in\mathcal{A}}\bar\beta_j^2
\E\{(\lambda_j^2)^{-1}\}+\E(\xi^{-1}),\\
b_\xi
&=\frac{1}{\tau_0^2}+\E\{(\tau^2)^{-1}\},\\
a_\zeta^\star
&=a_\zeta+\frac{|\mathcal{A}|}{2},\qquad
b_\zeta^\star=b_\zeta+\frac{1}{2}\sum_{j\in\mathcal{A}}\bar\beta_j^2.
\label{eq:rhs_ns_vb_updates}
\end{align}
The third line is shorthand for
$b_\tau=\frac{1}{2}\sum_{j\in\mathcal{A}}\bar\beta_j^2
\E\{(\lambda_j^2)^{-1}\}+\E(\xi^{-1})$, where the expectation of
$(\lambda_j^2)^{-1}$ is component-specific inside the summation.
As in MCMC, optional warmup scheduling for the $\tau^2,\xi$ block affects when
the coordinate is updated, not the coordinate target itself.

The following term-by-term ELBO contribution matches the package's
\code{rhs\_ns} scale block. For $x\sim\IG(a,b)$ under the rate
parameterization used above, define
\begin{align}
m_{\log x}(a,b)&=\log b-\psi(a),&
m_{x^{-1}}(a,b)&=a/b,\\
h_{\IG}(a,b)&=a+\log b+\log\Gamma(a)-(a+1)\psi(a),
\end{align}
where $\psi(\cdot)$ is the digamma function and $h_{\IG}$ is the entropy of an
inverse-gamma variational factor. Let
\begin{align}
\ell_{\lambda j}&=m_{\log x}(a_{\lambda j},b_{\lambda j}),&
r_{\lambda j}&=m_{x^{-1}}(a_{\lambda j},b_{\lambda j}),\\
\ell_{\nu j}&=m_{\log x}(a_{\nu j},b_{\nu j}),&
r_{\nu j}&=m_{x^{-1}}(a_{\nu j},b_{\nu j}),\\
\ell_\tau&=m_{\log x}(a_\tau,b_\tau),&
r_\tau&=m_{x^{-1}}(a_\tau,b_\tau),\\
\ell_\xi&=m_{\log x}(a_\xi,b_\xi),&
r_\xi&=m_{x^{-1}}(a_\xi,b_\xi).
\end{align}
If $\zeta^2$ is random, set
$\ell_\zeta=m_{\log x}(a_\zeta^\star,b_\zeta^\star)$ and
$r_\zeta=m_{x^{-1}}(a_\zeta^\star,b_\zeta^\star)$; if it is fixed at
$\zeta_0^2$, set $\ell_\zeta=\log\zeta_0^2$ and
$r_\zeta=1/\zeta_0^2$.

The coefficient-prior contribution is
\begin{align}
\elbo_{\beta,\mathrm{hs}}
&=
\sum_{j\in\mathcal{A}}
\left\{
-\frac12\log(2\pi)
-\frac12\ell_\tau
-\frac12\ell_{\lambda j}
-\frac12\bar\beta_j^2 r_\tau r_{\lambda j}
\right\},\\
\elbo_{\beta,\mathrm{slab}}
&=
\sum_{j\in\mathcal{A}}
\left\{
-\frac12\log(2\pi)
-\frac12\ell_\zeta
-\frac12\bar\beta_j^2 r_\zeta
\right\}.
\end{align}
The local, global, and slab prior contributions are
\begin{align}
\elbo_{\lambda\mid\nu}
&=
\sum_{j\in\mathcal{A}}
\left\{
-\frac12\ell_{\nu j}
-\log\Gamma(1/2)
-\frac32\ell_{\lambda j}
-r_{\nu j}r_{\lambda j}
\right\},\\
\elbo_{\nu}
&=
\sum_{j\in\mathcal{A}}
\left\{
-\log\Gamma(1/2)
-\frac32\ell_{\nu j}
-r_{\nu j}
\right\},\\
\elbo_{\tau\mid\xi}
&=
-\frac12\ell_\xi
-\log\Gamma(1/2)
-\frac32\ell_\tau
-r_\xi r_\tau,\\
\elbo_{\xi}
&=
\frac12\log(\tau_0^{-2})
-\log\Gamma(1/2)
-\frac32\ell_\xi
-\tau_0^{-2}r_\xi,\\
\elbo_{\zeta}
&=
\begin{cases}
a_\zeta\log b_\zeta-\log\Gamma(a_\zeta)
-(a_\zeta+1)\ell_\zeta-b_\zeta r_\zeta,
& \zeta^2\text{ random},\\
0, & \zeta^2\text{ fixed}.
\end{cases}
\end{align}
The entropy contribution is
\begin{equation}
\elbo_{\mathrm{ent}}
=
\sum_{j\in\mathcal{A}} h_{\IG}(a_{\lambda j},b_{\lambda j})
+\sum_{j\in\mathcal{A}} h_{\IG}(a_{\nu j},b_{\nu j})
+h_{\IG}(a_\tau,b_\tau)
+h_{\IG}(a_\xi,b_\xi)
+\Ind(\zeta^2\text{ random})h_{\IG}(a_\zeta^\star,b_\zeta^\star).
\end{equation}
Finally, when the intercept is not included in $\mathcal{A}$, the package adds
the fixed Gaussian intercept-prior contribution
\begin{equation}
\elbo_{\beta_0}
=
\frac12\{\log(\pi_0)-\log(2\pi)\}
-\frac12\pi_0\bar\beta_0^2,
\end{equation}
where $\pi_0$ is the fixed intercept precision. If the intercept is shrunk,
this term is omitted because the intercept is already represented in the
active \code{rhs\_ns} block.

Thus the package contribution is
\begin{equation}
\elbo_{\mathrm{rhs\_ns}}
=
\elbo_{\beta,\mathrm{hs}}
+\elbo_{\beta,\mathrm{slab}}
+\elbo_{\lambda\mid\nu}
+\elbo_{\nu}
+\elbo_{\tau\mid\xi}
+\elbo_{\xi}
+\elbo_{\zeta}
+\elbo_{\mathrm{ent}}
+\elbo_{\beta_0}.
\label{eq:rhs_ns_elbo}
\end{equation}
All terms are available in closed form because each scale factor is
inverse-gamma. This is the \code{rhs\_ns} component of the full static-model ELBO; the
likelihood, latent-variable, and $(\sigma,\gamma)$ terms are added by the
surrounding static AL or exAL variational routine.

\section{Forecasting, diagnostics, and posterior-predictive synthesis}
\label{sec:forecast_diagnostics_synthesis}

\subsection{Forecasting}

For a posterior draw $m$ at forecast origin $\tilde t$, the state forecast
recursion in the main article gives
\begin{equation}
\vect{\theta}_{\tilde t+k}^{(m)}
\sim
\N\{a_{\tilde t}^{(m)}(k),R_{\tilde t}^{(m)}(k)\},
\qquad k=1,\ldots,K.
\end{equation}
Given a forecasted state draw, posterior predictive draws are generated from
\begin{equation}
Y_{\tilde t+k}^{\mathrm{rep},(m)}
\sim
\exAL_{p_0}\{
\mat{F}_{\tilde t+k}^\top\vect{\theta}_{\tilde t+k}^{(m)},
\sigma^{(m)},\gamma^{(m)}
\}.
\label{eq:forecast_postpred}
\end{equation}
For the AL/DQLM special case, set $\gamma^{(m)}=0$. The function
\code{exdqlmForecast()} implements this recursion for MCMC and LDVB fitted
objects; for LDVB fits, the posterior draws are replaced by draws from the
fitted variational factors.

\subsection{Diagnostics}

The dynamic diagnostics use the plug-in maximum a posteriori (MAP) standardized one-step-ahead
forecast-error sequence stored as \code{map.standard.forecast.errors}. If
$\widehat f_t$ and $\widehat q_t$ denote the fitted one-step-ahead location and
variance used in this sequence, the diagnostic error and its normal-score
transform are
\begin{equation}
\widehat e_t = \frac{y_t-\widehat f_t}{\sqrt{\widehat q_t}},
\qquad
\widehat u_t = \Phi(\widehat e_t).
\label{eq:pit_plugin}
\end{equation}
The quantile--quantile (QQ) plot, autocorrelation function (ACF) plot, and standardized forecast-error plot are computed from
$\{\widehat e_t\}$ or $\{\widehat u_t\}$ rather than from posterior uncertainty
bands for residuals. The reported Kullback--Leibler (KL) summaries are nearest-neighbor estimates
of the divergence between this empirical standardized-error sample and a
reference standard-normal sample supplied through \code{ref} or generated
internally; the flipped version reverses the two arguments.

For predictive draws $y_t^{(1)},\ldots,y_t^{(M)}$, let
\(\widehat q_t(\tau_k)\) denote the empirical posterior predictive quantile at
level \(\tau_k\). The diagnostic continuous ranked probability score (CRPS)
calculation uses the identity between CRPS and integrated check loss,
\[
\CRPS(y_t^{\mathrm{obs}},G_t)
=
2\int_0^1
\rho_\tau\{y_t^{\mathrm{obs}}-G_t^{-1}(\tau)\}\,d\tau,
\]
and reports the finite integrated quantile-score approximation
\begin{equation}
\widehat{\CRPS}_t
=
2\sum_{k=1}^{K}
w_k
\rho_{\tau_k}\{y_t^{\mathrm{obs}}-\widehat q_t(\tau_k)\}.
\label{eq:crps_estimator}
\end{equation}
The reported CRPS is the average of \eqref{eq:crps_estimator} over the
diagnostic time points \citep{gneiting2007,laio2007verification}. By default,
\code{exdqlmDiagnostics()} uses
\(\tau_k=0.01,0.02,\ldots,0.99\) with equal weights \(w_k=1/K\). User-supplied
weights are normalized to sum to one. For target quantile $p_0$, the
posterior predictive loss criterion (PPLC) returned as \code{pplc} is
\begin{equation}
\widehat{\PPLC}
=
\sum_t
\frac{1}{M}\sum_{m=1}^M
\rho_{p_0}\{y_t^{\mathrm{obs}}-y_t^{\mathrm{rep},(m)}\},
\label{eq:pplc_estimator}
\end{equation}
where $\rho_{p_0}$ is the check-loss function \citep{gelfand1998model}.
The dynamic diagnostic routine returns these summaries. The static diagnostic
routine returns check-loss summaries and reference-quantile errors when a
reference quantile is supplied.

\subsection{Posterior-predictive synthesis}

The function \code{quantileSynthesis()} combines posterior predictive draws
from separately fitted models at ordered levels $p_1<\cdots<p_K$. For each time
point and draw index, it forms the quantile-specific predictive curve, applies
the requested monotonicity correction when fitted quantiles cross, and
interpolates across $p$ to obtain draws from one posterior predictive
distribution. The monotonicity correction uses isotonic adjustment and, when
requested, global monotone rearrangement \citep{barlow1972,chernozhukov2010}.

\begin{suppalgorithm}{Posterior-predictive synthesis}
\label{alg:predictive-synthesis}
\par\noindent\emph{Input:} posterior predictive draws or fitted forecast
objects from models at ordered quantile levels $p_1<\cdots<p_K$, together with
the requested monotonicity correction and interpolation settings.
\par\noindent\emph{Output:} synthesized posterior predictive draws and
summary intervals on the common predictive scale.
\begin{enumerate}[label=(\arabic*)]
\item Extract or collect posterior predictive draw matrices for each fitted
quantile level.
\item For each time point and retained draw index, form the quantile-specific
predictive curve over $p_1,\ldots,p_K$.
\item Apply the requested monotonicity correction.
\item Interpolate across quantile levels.
\item Return synthesized posterior predictive draws and summary intervals.
\end{enumerate}
\end{suppalgorithm}

\section{Numerical and ELBO implementation blocks}
\label{sec:numerical_blocks}

\subsection{Reusable moment and entropy formulas}
\label{sec:elbo_building_blocks}

The variational routines monitor
\begin{equation}
\elbo(q)=\E_q\{\log p(\text{complete data})\}+H(q),
\qquad
H(q)=-\E_q\{\log q\}.
\end{equation}
The following moment and entropy formulas are reused by the dynamic and static
algorithms.

If $\vect{x}\sim\N(\vect{m},\mat{V})$ with
$\vect{x}\in\R^d$, then
\begin{equation}
H\{\N(\vect{m},\mat{V})\}
=\frac{1}{2}\{d(1+\log 2\pi)+\log|\mat{V}|\}.
\label{eq:entropy_gaussian}
\end{equation}
If $x\sim\IG(a,b)$ under the parameterization in
Section~\ref{sec:supp_distributional_conventions}, then
\begin{equation}
\E(x^{-1})=\frac{a}{b},\qquad
\E(\log x)=\log b-\psi(a),
\end{equation}
and
\begin{equation}
H\{\IG(a,b)\}
=a+\log b+\log\Gamma(a)-(a+1)\psi(a),
\label{eq:entropy_ig}
\end{equation}
where $\psi(\cdot)$ is the digamma function.

If $x\sim\GIG(\lambda,\chi,\psi_{\mathrm{GIG}})$ and
$z=(\chi\psi_{\mathrm{GIG}})^{1/2}$, let
$m_x=\E(x)$, $m_{1/x}=\E(x^{-1})$, and
$m_{\log x}=\E(\log x)$. Then
\begin{align}
H\{\GIG(\lambda,\chi,\psi_{\mathrm{GIG}})\}
&=
\log\{2K_\lambda(z)\}
-\frac{\lambda}{2}\log\left(\frac{\psi_{\mathrm{GIG}}}{\chi}\right)
-(\lambda-1)m_{\log x} \nonumber\\
&\quad+
\frac{1}{2}\{\chi m_{1/x}+\psi_{\mathrm{GIG}} m_x\}.
\label{eq:entropy_gig}
\end{align}
The moments $m_x$ and $m_{1/x}$ follow from \eqref{eq:gig_moments};
$m_{\log x}$ is the derivative of $\log K_\lambda(z)$ with respect to
$\lambda$ plus $\frac{1}{2}\log(\chi/\psi_{\mathrm{GIG}})$, and is evaluated
numerically in the package.

If $x\sim\TNp(m,V)$, let $a=V^{1/2}$, $z=m/a$, and
$\Lambda(z)=\phi(z)/\Phi(z)$. The entropy is
\begin{equation}
H\{\TNp(m,V)\}
=
\frac{1}{2}\log(2\pi V)+\log\Phi(z)
+\frac{1}{2}\{1-z\Lambda(z)\}.
\label{eq:entropy_tnp}
\end{equation}
This closed form is used for the truncated-normal latent factors; numerically,
the package evaluates the required moments using stable tail-probability
guards.

For exAL variational inference, the nonconjugate $(\sigma,\gamma)$ block is handled on
transformed coordinates
\begin{equation}
\ell_\sigma=\log\sigma,\qquad
z_\gamma=\logit\left(\frac{\gamma-L}{U-L}\right).
\end{equation}
If $\vect{z}=(z_\gamma,\ell_\sigma)^\top$ has Laplace approximation
$q_z(\vect{z})=\N(\hat{\vect{z}},\mat{\Sigma}_z)$ and
$(\gamma,\sigma)=T(\vect{z})$, then the entropy contribution on the original
scale is approximated as
\begin{equation}
H\{q(\sigma,\gamma)\}
\approx
\frac{1}{2}\{2(1+\log 2\pi)+\log|\mat{\Sigma}_z|\}
+\E_{q_z}\{\log|J_T(\vect{z})|\},
\label{eq:entropy_ld}
\end{equation}
where
\begin{equation}
\log|J_T(\vect{z})|
=\ell_\sigma+\log(U-L)+\log r+\log(1-r),
\qquad r=\expit(z_\gamma).
\end{equation}
The final expectation in \eqref{eq:entropy_ld}, and other smooth expectations
of functions of $(\sigma,\gamma)$, are evaluated by the same second-order
delta approximation used in the LDVB updates.

\subsection{Bounded univariate slice sampler}
\label{sec:slice}

Let $h(x)$ be an unnormalized log-density on a finite interval $[a,b]$, and let
$x_0\in(a,b)$. Following the stepping-out and shrinkage construction of
\citet{neal2003slice}, the bounded slice update implemented in the package is:
\begin{enumerate}[label=(\arabic*)]
\item Draw $e\sim\Exp(1)$ and set the log-slice height $z=h(x_0)-e$.
\item Draw $u\sim\operatorname{Unif}(0,w)$ and initialize
$L=x_0-u$, $R=x_0+w-u$.
\item Step out left and right by increments of $w$, respecting the bounds
$a$ and $b$, until the interval endpoints lie below the slice or the maximum
number of expansions is reached.
\item Truncate the interval to $[a,b]$.
\item Repeatedly draw $x_1\sim\operatorname{Unif}(L,R)$. If $h(x_1)\ge z$,
accept $x_1$. Otherwise shrink the interval to $[L,x_1]$ if $x_1>x_0$ and to
$[x_1,R]$ otherwise.
\end{enumerate}
For dynamic and static exAL MCMC fits, this sampler is applied to the
one-dimensional $\gamma$ target on $(L,U)$ by default.

\subsection{Laplace--delta approximation}

For a nonconjugate block $\vect{\eta}\in\R^d$ with log target
$\ell(\vect{\eta})$, let $\hat{\vect{\eta}}$ be a local mode and
$\mat{H}=\nabla^2\ell(\hat{\vect{\eta}})$. If $-\mat{H}$ is positive definite,
the Laplace approximation is
\begin{equation}
q(\vect{\eta})\approx
\N\left(\hat{\vect{\eta}},\ [-\mat{H}]^{-1}\right).
\label{eq:ld_gaussian}
\end{equation}
For a smooth function $r(\vect{\eta})$, the delta approximation used by LDVB is
\begin{equation}
\E\{r(\vect{\eta})\}
\approx
r(\hat{\vect{\eta}})
+
\frac{1}{2}\tr\left[
\nabla^2 r(\hat{\vect{\eta}})[-\mat{H}]^{-1}
\right].
\label{eq:delta_generic}
\end{equation}
This is used for smooth functions of $(\sigma,\gamma)$ in dynamic and static
exAL models.

\clearpage
\bibliography{references}

\end{document}
